\documentclass[11pt]{article}
\usepackage[T1]{fontenc}
\usepackage{lmodern}
\usepackage[utf8]{inputenc}
\usepackage{microtype}
\usepackage[a4paper,margin=27mm,headheight=15pt]{geometry}
\usepackage{amsmath,amssymb,amsthm,mathtools}
\usepackage{booktabs,longtable,array,tabularx}
\usepackage{enumitem}
\usepackage{xcolor}
\usepackage{hyperref}
\usepackage[nameinlink,noabbrev]{cleveref}
\usepackage{fancyhdr}
\usepackage{listings}
\usepackage{verbatim}

\hypersetup{
  colorlinks=true,
  linkcolor=blue!45!black,
  citecolor=green!35!black,
  urlcolor=blue!55!black,
  pdftitle={Expected value in intuitionistic primitive recursive arithmetic},
  pdfauthor={OpenAI, prepared for Vladimir Reshetnikov}
}

\pagestyle{fancy}
\fancyhf{}
\fancyhead[L]{Expected value in iPRA}
\fancyhead[R]{16 November 2026}
\fancyfoot[C]{\thepage}

\newtheorem{theorem}{Theorem}[section]
\newtheorem{lemma}[theorem]{Lemma}
\newtheorem{proposition}[theorem]{Proposition}
\newtheorem{corollary}[theorem]{Corollary}
\newtheorem{definition}[theorem]{Definition}
\newtheorem{example}[theorem]{Example}
\newtheorem{remark}[theorem]{Remark}
\newtheorem{construction}[theorem]{Construction}
\newtheorem{contract}[theorem]{Effective contract}
\newtheorem{audit}[theorem]{Audit}

\newcommand{\N}{\mathbb N}
\newcommand{\Z}{\mathbb Z}
\newcommand{\Q}{\mathbb Q}
\newcommand{\R}{\mathbb R}
\newcommand{\E}{\mathbb E}
\newcommand{\Pp}{\mathbb P}
\newcommand{\Var}{\operatorname{Var}}
\newcommand{\Cov}{\operatorname{Cov}}
\newcommand{\ind}{\mathbf 1}
\newcommand{\abs}[1]{\lvert #1\rvert}
\newcommand{\norm}[1]{\lVert #1\rVert}
\newcommand{\code}[1]{\texttt{\detokenize{#1}}}
\newcommand{\ipra}{\mathrm{iPRA}}
\newcommand{\qfind}{\mathrm{QF\mbox{-}Ind}}
\newcommand{\Fin}{\operatorname{Fin}}
\newcommand{\supp}{\operatorname{supp}}
\newcommand{\sgn}{\operatorname{sgn}}
\newcommand{\liminfN}{\mathop{\mathrm{liminf}}}
\newcommand{\limsupN}{\mathop{\mathrm{limsup}}}
\newcommand{\defeq}{\mathrel{:=}}
\newcommand{\dyad}[1]{2^{-#1}}

\setlist[itemize]{leftmargin=2em,itemsep=2pt,topsep=4pt}
\setlist[enumerate]{leftmargin=2.2em,itemsep=2pt,topsep=4pt}
\lstset{basicstyle=\ttfamily\small,breaklines=true,columns=fullflexible,frame=single}

\title{Expected Value in Intuitionistic Primitive Recursive Arithmetic\\
\large Finite exact probability, effective limits, and a Lean formalization route}
\author{OpenAI\\Prepared for Vladimir Reshetnikov}
\date{16 November 2026}

\begin{document}
\maketitle

\begin{abstract}
This report reconstructs as much as possible of the mathematical content of the
Wikipedia article \emph{Expected value} inside intuitionistic primitive
recursive arithmetic, abbreviated \(\ipra\).  The central observation is that a
large majority of the article's algebraic probability theory is already
finitary.  On a finite sample space with natural weights and rational-valued
observables, expectation is one bounded sum divided by a positive integer.
Normalization, linearity, monotonicity, the law of the unconscious statistician,
indicator and tail formulas, variance identities, Markov and Chebyshev
inequalities, Cauchy--Schwarz, independence, finite conditional expectation,
and sample-mean formulas reduce to primitive-recursive calculations and
quantifier-free induction.

The article also discusses countable laws, densities, general Lebesgue
integration, convergence theorems, characteristic functions, and applications.
Those statements are not uniformly meaningful in a theory whose only objects
are natural numbers.  We give precise effective replacements: every infinite
object is a primitive-recursive stream of finite rational approximants together
with an explicit tail, Cauchy, integrability, or continuity modulus.  The
resulting claims are theorem schemata of \(\ipra\), not appeals to an ambient
classical universe.  We isolate where a classical converse---most notably an
existence-only Radon--Nikodym assertion---would require additional data and is
therefore not claimed.

The manuscript supplies complete human-readable arguments, a theorem-by-theorem
coverage ledger, three Lean design files, a route to explicit derivations in
ProveIt's \code{Prf PRAAx}, and an exact-rational regression program.  The Lean
files are blueprints rather than a completed kernel-checked formalization.  The
Python run checks 148,317 exact assertions and uses no floating-point value as a
proof oracle.
\end{abstract}

\vspace{1em}
\noindent\textbf{Formal status.}
The mathematical proofs below are written so that every finite step can be
compiled to bounded sums, products, comparisons, pairing, and
quantifier-free induction.  The effective analytic arguments state all moduli
needed for such a compilation.  Nevertheless, this packet does not contain the
final explicit \code{PRAAx} proof terms.  Statements marked ``internal'' are
intended iPRA theorems; statements marked ``effective analytic'' are iPRA
schemata for a fixed primitive-recursive presentation and its certificates.

\tableofcontents
\newpage

\section*{Executive theorem map}
\addcontentsline{toc}{section}{Executive theorem map}

The paper has three layers.

\begin{enumerate}
\item \textbf{Exact finite layer.}  \Cref{sec:coding,sec:finite,sec:properties,sec:inequalities,sec:independence,sec:samples}
  develops expectation on a finite weighted law.  No real numbers, completed
  infinite sets, or measure spaces are used.  After clearing denominators, the
  statements are equations and inequalities over \(\N\).
\item \textbf{Effective limiting layer.}  \Cref{sec:countable,sec:integration,sec:distributions,sec:convergence}
  treats series and integrals through finite approximants and explicit moduli.
  The finite kernel is applied at a precision computed from the requested output
  precision.
\item \textbf{Formalization layer.}  \Cref{sec:fourier,sec:lean,sec:conclusion}
  covers finite Fourier methods and applications, records the Lean interfaces,
  and identifies the remaining proof-engineering work.
\end{enumerate}

\begin{center}
\begin{tabular}{@{}p{0.24\textwidth}p{0.66\textwidth}@{}}
\toprule
Class & Representative results \\
\midrule
Exact in \(\ipra\) & finite expectation and LOTUS; linearity; monotonicity;
positive/negative parts; indicator and tail formulas; variance; Markov,
Chebyshev and Cauchy--Schwarz; product-law independence; finite conditioning;
sample mean; finite Kolmogorov inequality. \\
Effective schema & absolutely convergent countable expectation; effective
reordering; lower-real extended expectation; compact and improper integration;
monotone/Fatou/uniform-integrability convergence; common continuous laws;
characteristic-function differentiation. \\
Constructive correction & conditionally convergent sums are order sensitive;
Cauchy truncations are path dependent; pointwise convergence does not imply
convergence of expectations; the article's ``about 53 percent'' die statement
corresponds to an open interval, not the inclusive event. \\
Not claimed uniformly & extraction of a density from bare absolute continuity;
unwitnessed arbitrary Lebesgue objects; qualitative almost-sure convergence
without a finite-error contract or a modulus presentation. \\
\bottomrule
\end{tabular}
\end{center}

\section{Constructive scope and reading conventions}\label{sec:scope}

\subsection{The base theory}

We use \(\ipra\) in the standard first-order sense: the only individual sort is
\(\N\); every primitive-recursive operation may be named by a function symbol;
its defining equations are available; reasoning is intuitionistic; and
induction is restricted to quantifier-free formulas.  Finite sequences,
finite trees, signed integers, rationals, and tuples are therefore represented
by natural-number codes.  Bounded search and bounded sums are primitive
recursive.  A theorem about a coded finite object is internal when its proof is
an explicit derivation from those defining equations and quantifier-free
induction.

The restriction to quantifier-free induction is not an obstacle to the finite
probability identities in the article.  Each relevant invariant can be
maintained by an equation or decidable inequality.  For example, the statement
that the first \(n\) terms of two finite sums agree is an equation between two
primitive-recursive terms.  Its successor step is one use of the defining
recurrence for bounded summation.

\begin{remark}[No hidden classical quotient]
A rational number is not treated as an abstract quotient field.  It is a code
with a positive denominator, and equalities are proved by
cross-multiplication.  Likewise, ``almost surely'' on a finite weighted law is
an equation saying that the total weight of the exceptional outcomes is zero.
Both are decidable arithmetical relations.
\end{remark}

\subsection{Three meanings of a theorem}

The article mixes finite algebra, countable series, and general integration.
We use the following disciplined translations.

\begin{definition}[Exact finite theorem]
An exact finite theorem quantifies over natural codes for a finite law and
finite observables.  After denominator clearing, its conclusion is a
quantifier-free equation or inequality.  The proof uses only primitive
recursion and quantifier-free induction.
\end{definition}

\begin{definition}[Effective analytic theorem schema]
An effective analytic theorem is a uniform family of finite theorems indexed by
an output precision \(p\).  The input contains primitive-recursive approximants
and explicit moduli.  From \(p\), the proof computes a finite stage \(N(p)\),
applies exact finite reasoning there, and proves an error at most \(2^{-p}\).
For each fixed presentation, the entire checker is primitive recursive.
\end{definition}

\begin{definition}[Constructive replacement]
When the article states a classical existence theorem whose data do not contain
an algorithm or modulus, we replace it by a witnessed version.  We prove the
forward implications from the witness and record that the non-effective
converse has not been internalized.
\end{definition}

\subsection{What ``prove as many statements as possible'' means}

The theorem ledger accompanying this paper separates the article's content
into mathematical claims, examples, and historical prose.  We prove every
finite algebraic claim, including the standard inequalities and product-law
formulas.  We also prove effective versions of the main countable and
continuous claims.  Some article sentences are definitions rather than
propositions; in those cases the task is to show that the definition has a
primitive-recursive realization and interacts correctly with the finite
kernel.

We do not claim that every classical random variable or measure has an
\(\ipra\)-usable name.  A bare statement such as ``let \(X\) be an integrable
random variable'' suppresses both a representation and a convergence rate.
Our version instead says: let \(X_p\) be a primitive-recursive sequence of
finite rational simple laws with a certified \(L^1\) Cauchy bound.  This is
strictly more informative and is exactly what a proof checker can consume.

\subsection{Constructive logic in the proofs}

All finite predicates occurring below are decidable: equality and order of
coded integers and rationals, membership in a finite event, and comparison of a
bounded sum with a rational threshold.  Case distinctions on such predicates
are intuitionistically valid.  We never infer an arbitrary proposition \(A\)
from \(\neg\neg A\), and we never use an unbounded least counterexample.

Several classical-looking arguments become simpler when written pointwise.
Markov's inequality is the pointwise inequality
\(a\ind_{X\ge a}\le X\), summed with nonnegative weights.  The variance identity
is a polynomial identity.  Cauchy--Schwarz is a sum of squares.  The law of
total expectation is a finite regrouping of one double sum.  These proofs do
not require excluded middle beyond decidable comparisons already supplied by
the codes.

\subsection{Trust boundary}

The trusted mathematical boundary intended by this report is the existing
ProveIt calculus \code{Prf PRAAx} and Lean's kernel.  The supplied Python script
is deliberately outside that boundary: it exhaustively and randomly checks
many small exact instances, but it is only regression evidence.  Similarly,
ambient Lean structures in the packet are specifications until their fields
are filled by actual proofs and then translated to object-language derivations.

\section{Coding finite probability in iPRA}\label{sec:coding}

\subsection{Signed integers and rationals}

A signed integer is represented by a pair \((a,b)\in\N^2\), denoting
\(a-b\).  Equality and order are
\[
 (a,b)=(c,d) \quad\Longleftrightarrow\quad a+d=c+b,
\]
\[
 (a,b)\le(c,d) \quad\Longleftrightarrow\quad a+d\le c+b.
\]
Addition, multiplication, negation, and absolute value are primitive recursive.
For instance,
\((a,b)(c,d)=(ac+bd,ad+bc)\), and a decidable comparison selects the positive
or negative part for absolute value.

A rational code is \((z,q)\), where \(z\) is a signed-integer code and
\(q>0\).  It denotes \(z/q\).  Equality and order are again decided by positive
cross-multiplication.  Reduction to lowest terms is optional; no proof below
relies on unique representation.

\begin{lemma}[Common denominator normalization]\label{lem:common-den}
For every coded finite list of rational numbers
\(x_0,\ldots,x_{m-1}\), there is a primitive-recursively computed positive
integer \(D\) and signed integers \(a_i\) such that \(x_i=a_i/D\) for every
\(i<m\).
\end{lemma}

\begin{proof}
Take the product of the positive input denominators.  Multiplying each
numerator by the product of the remaining denominators gives the required
\(a_i\).  The bounded product and the transformed list are primitive recursive.
All displayed equalities follow by distributivity; no gcd computation is
needed.
\end{proof}

\subsection{Finite weighted laws}

\begin{definition}[Finite law]\label{def:finlaw}
A finite weighted law is coded by
\[
 m\in\N,\qquad w:\Fin(m)\to\N,
 \qquad W=\sum_{i<m}w_i>0.
\]
The probability of an event \(A\subseteq\Fin(m)\), represented by a decidable
indicator \(a_i\in\{0,1\}\), is
\[
 \Pp(A)=\frac{1}{W}\sum_{i<m}w_i a_i.
\]
A rational observable is a finite list \(X_i\in\Q\), and its expectation is
\[
 \E[X]=\frac{1}{W}\sum_{i<m}w_iX_i.
\]
\end{definition}

The use of natural weights loses no finite rational law.  If the listed
probabilities are \(p_i=r_i/s_i\), choose a common denominator \(W\), put
\(w_i=Wp_i\), and delete zero-weight padding if desired.  Conversely, natural
weights produce the rational probabilities \(w_i/W\).

\begin{lemma}[Primitive-recursive evaluator]\label{lem:pr-evaluator}
Given codes for a finite law and a rational observable, a code for
\(\E[X]\) is primitive recursive.  Equality and order between two such
expectations are quantifier-free after clearing the positive denominators.
\end{lemma}

\begin{proof}
Use \Cref{lem:common-den} to write \(X_i=a_i/D\).  Then
\[
 \E[X]=\frac{\sum_{i<m}w_i a_i}{WD}.
\]
The numerator is a bounded sum of primitive-recursive signed products, and
\(WD>0\).  Comparing two expectations is therefore a comparison of two signed
cross-products.
\end{proof}

\subsection{Events, null sets, and almost-sure equality}

\begin{definition}[Finite almost-sure equality]\label{def:ae}
For observables \(X,Y\) on the same finite law, write \(X=Y\) almost surely
when
\[
  \sum_{i<m} w_i\ind_{X_i\ne Y_i}=0.
\]
An event is null when its weighted indicator sum is zero.
\end{definition}

This definition is stronger computationally than merely saying that the
probability quotient equals zero: its numerator is explicitly zero.  Since all
summands are natural and nonnegative, zero-sum elimination gives
\(w_i\ind_{X_i\ne Y_i}=0\) for every \(i<m\).  ProveIt already contains exactly
this bounded-sum pattern: a zero bounded sum has zero terms, and conversely
pointwise zero terms give a zero sum.

\subsection{Compiling bounded algebra to quantifier-free induction}

\begin{proposition}[Finite-sum compilation pattern]\label{prop:sum-compile}
Each of the following has an explicit \(\ipra\) derivation by
quantifier-free induction on the bound:
\begin{align*}
 \sum_{i<n+1}f(i)&=\sum_{i<n}f(i)+f(n),\\
 \sum_{i<n}(f(i)+g(i))&=\sum_{i<n}f(i)+\sum_{i<n}g(i),\\
 \left[\forall i<n\;f(i)=g(i)\right]&\Longrightarrow
     \sum_{i<n}f(i)=\sum_{i<n}g(i),\\
 j<n&\Longrightarrow f(j)\le\sum_{i<n}f(i)
 \quad\text{when all values are natural.}
\end{align*}
The bounded universal premise is represented by a zero sum of decidable
failure indicators.
\end{proposition}

\begin{proof}
The first equation is the defining equation of the primitive-recursive sum.
For the second, the induction invariant is the displayed equality.  The base
case is \(0=0+0\); the successor step substitutes the first recurrence and
uses associativity and commutativity.  For congruence, maintain equality of the
partial sums and use the zero failure-count to obtain equality at the new
index.  The final inequality follows because the target term occurs as one
nonnegative summand; an induction on the upper bound gives the formal proof.
All invariants are quantifier-free equations or inequalities.
\end{proof}

\begin{corollary}[Finite Fubini and regrouping]\label{cor:finite-fubini}
For coded finite arrays,
\[
 \sum_{i<m}\sum_{j<n}a_{ij}
 =\sum_{j<n}\sum_{i<m}a_{ij}
 =\sum_{k<mn}a_{\lfloor k/n\rfloor,\,k\bmod n}.
\]
More generally, a finite sum may be regrouped by any primitive-recursive map
between finite index sets, provided multiplicity is tracked.
\end{corollary}

\begin{proof}
Induct on one rectangle dimension and use distributivity.  The mixed-radix
map \(k\mapsto(\lfloor k/n\rfloor,k\bmod n)\) is primitive recursive by division
with remainder.  The two directions of the bijection are verified by the
usual quotient-remainder equations, which are available in the arithmetic
library.
\end{proof}

\subsection{Rational inequalities without division}

Whenever a theorem contains fractions, the formal target is its
positive-denominator clearing.  For example, monotonicity
\[
 \frac{A}{WD}\le\frac{B}{WE}
\]
becomes \(AE\le BD\) after checking the correct denominators and signs.  A
proof should never invoke an abstract ordered field when a direct natural
identity is available.  This policy both shortens the object-language proof
and makes the computational content visible.

\section{The exact finite expectation kernel}\label{sec:finite}

Throughout this section \(L=(m,w,W)\) is a finite law and
\(X,Y:\Fin(m)\to\Q\) are rational observables.  Every identity is understood
both as ordinary rational mathematics and as its denominator-cleared coded
form.

\begin{definition}[Expectation, moment, and event probability]\label{def:expect}
Define
\[
 \E_L[X]\defeq\frac{\sum_{i<m}w_iX_i}{W},\qquad
 \E_L[X^r]\defeq\frac{\sum_{i<m}w_iX_i^r}{W},
\]
and, for a decidable event \(A\),
\[
 \Pp_L(A)\defeq\E_L[\ind_A]
 =\frac{\sum_{i<m}w_i\ind_A(i)}{W}.
\]
\end{definition}

\begin{theorem}[Normalization]\label{thm:normalization}
\(\Pp_L(\Fin(m))=1\), \(\Pp_L(\varnothing)=0\), and
\(0\le\Pp_L(A)\le1\) for every event \(A\).
\end{theorem}

\begin{proof}
For the whole space the indicator is constantly one, so its weighted sum is
\(\sum_iw_i=W\).  For the empty event it is constantly zero.  Finally,
\(0\le w_i\ind_A(i)\le w_i\) pointwise; bounded-sum monotonicity gives
\(0\le\sum_iw_i\ind_A(i)\le W\), and division by the positive \(W\) preserves
order.  Internally one retains the cross-multiplied numerator inequalities.
\end{proof}

\begin{theorem}[Weighted-average formula]\label{thm:weighted-average}
If \(p_i=w_i/W\), then \(p_i\ge0\), \(\sum_i p_i=1\), and
\[
 \E_L[X]=\sum_{i<m}p_iX_i.
\]
Conversely every finite rational probability vector has such a natural-weight
presentation.
\end{theorem}

\begin{proof}
The forward statements are direct expansion.  For the converse, take the
product of all positive denominators as a common \(W\) and put \(w_i=Wp_i\).
The equation \(\sum p_i=1\) becomes \(\sum w_i=W\) after multiplication by
\(W\).
\end{proof}

\begin{theorem}[Convex-hull bounds]\label{thm:convex-hull}
Suppose rational numbers \(a,b\) satisfy
\(a\le X_i\le b\) for every positive-weight outcome.  Then
\[
 a\le\E[X]\le b.
\]
If \(\E[X]=a\), then \(X_i=a\) at every positive-weight outcome whenever
\(X_i\ge a\); similarly at the upper endpoint.
\end{theorem}

\begin{proof}
Pointwise, \(w_i a\le w_iX_i\le w_i b\).  Sum and use
\(\sum_iw_i=W\).  For the endpoint assertion write \(X_i-a\ge0\).  The equality
\(\E[X]=a\) says \(\sum_iw_i(X_i-a)=0\) after denominator clearing.  Every
nonnegative summand is therefore zero.  If \(w_i>0\), cancellation yields
\(X_i=a\).
\end{proof}

\begin{corollary}[Equiprobable arithmetic mean]\label{cor:equiprobable}
For equal positive weights,
\[
 \E[X]=\frac{X_0+\cdots+X_{m-1}}{m}.
\]
In particular, expectation generalizes the ordinary arithmetic mean.
\end{corollary}

\begin{proof}
Take \(w_i=1\) and \(W=m\) in \Cref{def:expect}.
\end{proof}

\subsection{Pushforwards and finite LOTUS}

\begin{definition}[Pushforward law]\label{def:pushforward}
Let \(f:\Fin(m)\to\Fin(n)\) be primitive recursive.  Its pushforward weights
are
\[
 v_j\defeq\sum_{\substack{i<m\\f(i)=j}}w_i.
\]
The total remains \(\sum_{j<n}v_j=W\).
\end{definition}

\begin{theorem}[Finite law of the unconscious statistician]\label{thm:finite-lotus}
For every rational \(g:\Fin(n)\to\Q\),
\[
 \E_L[g(f(i))]=\frac1W\sum_{j<n}v_jg(j).
\]
Equivalently, expectation after applying a function depends only on the
pushforward distribution.
\end{theorem}

\begin{proof}
Insert a decidable partition of unity:
\[
 \sum_{i<m}w_i g(f(i))
 =\sum_{i<m}\sum_{j<n}w_i\ind_{f(i)=j}g(j).
\]
For each \(i\), exactly one indicator is one, so the inner sum equals the
original term.  Apply finite Fubini and factor out \(g(j)\):
\[
 \sum_{j<n}g(j)\sum_{i<m}w_i\ind_{f(i)=j}
 =\sum_{j<n}v_jg(j).
\]
Every equality is a bounded-sum identity.  The claim that exactly one fiber
indicator is one is proved by decidable equality on \(\Fin(n)\).
\end{proof}

\begin{corollary}[Multidimensional LOTUS]\label{cor:multi-lotus}
For finite observables \(X_1,\ldots,X_r\) and a rational function \(g\) on
their finite joint range,
\[
 \E[g(X_1,\ldots,X_r)]
 =\sum_{x_1,\ldots,x_r}g(x_1,\ldots,x_r)
   \Pp(X_1=x_1,\ldots,X_r=x_r).
\]
\end{corollary}

\begin{proof}
Code the tuple by primitive-recursive pairing and apply
\Cref{thm:finite-lotus}.  Repeated finite Fubini converts the single paired
sum to the displayed nested sum.
\end{proof}

\subsection{Vector and matrix expectation}

\begin{theorem}[Componentwise expectation]\label{thm:componentwise}
Let \(X(i)=(X_1(i),\ldots,X_r(i))\in\Q^r\).  Define expectation of a vector by
its finite weighted sum.  Then
\[
 \E[X]=(\E[X_1],\ldots,\E[X_r]).
\]
For a rational matrix-valued observable, expectation is entrywise.  Rational
matrix multiplication by a fixed matrix commutes with expectation.
\end{theorem}

\begin{proof}
A coded vector or matrix is a finite tuple of rational codes.  Projection onto
a coordinate commutes definitionally with the bounded weighted sum.  For a
fixed matrix \(A\), each coordinate of \(AX\) is a finite linear combination
of coordinates of \(X\); distributivity of bounded sums proves
\(\E[AX]=A\E[X]\).  No higher-type object is needed: the dimension and all
coordinate access functions are natural-number parameters.
\end{proof}

\subsection{Finite distributions as normal forms}

\begin{proposition}[Histogram normal form]\label{prop:histogram}
Every finite rational observable has a primitive-recursively computed finite
histogram \((x_j,v_j)_{j<n}\) such that the \(x_j\) are pairwise distinct,
\(v_j\) is the total weight of outcomes with value \(x_j\), and
\[
 \E[X]=\frac1W\sum_{j<n}v_jx_j.
\]
The histogram may be obtained without sorting by retaining one slot per input
outcome and allowing zero duplicate slots.
\end{proposition}

\begin{proof}
Use the pushforward map from outcome indices to the least earlier index with
the same rational value; bounded search computes this representative.
Alternatively, define
\(v_j=w_j\ind_{j\text{ is first}}+\sum_{i>j}w_i\ind_{X_i=X_j}\ind_{j\text{ is first}}\).
Both are primitive-recursive.  Fiber regrouping gives the expectation identity.
The zero-padding option avoids any need for a variable-length quotient type.
\end{proof}

\begin{remark}
The histogram form is useful when translating textbook notation
\(\sum_x x\Pp(X=x)\).  The sample-space form is usually simpler for formal
proofs, because no proof of distinctness or enumeration completeness is
required.
\end{remark}

\section{Finite examples from the article}\label{sec:examples}

\begin{example}[Fair die]\label{ex:die}
For a fair six-sided die with outcomes \(1,2,3,4,5,6\),
\[
 \E[X]=\frac{1+2+3+4+5+6}{6}=\frac{21}{6}=\frac72.
\]
The formal proof is the natural identity \(2\cdot21=7\cdot6\), after evaluating
six successor recurrences of the bounded sum.
\end{example}

\begin{example}[American roulette straight-up bet]\label{ex:roulette}
A one-dollar straight-up bet on an American roulette wheel wins \(35\) dollars
of net profit on one of \(38\) slots and loses \(1\) dollar on the other
\(37\).  Hence
\[
 \E[G]=\frac{35-37}{38}=-\frac{2}{38}=-\frac1{19}.
\]
For \(190\) bets, whether or not the bets are independent, linearity gives
\[
 \E\left[\sum_{j<190}G_j\right]
 =190\left(-\frac1{19}\right)=-10.
\]
Independence matters for variance and concentration, not for this expectation
identity.
\end{example}

\begin{example}[Expectation need not be an attainable value]\label{ex:not-attained}
The die expectation \(7/2\) is not one of the possible die outcomes.  Thus the
word ``expected'' does not mean ``a value that must occur.''  In the finite
kernel, expectation is a rational barycenter of the support and is guaranteed
only to lie between its minimum and maximum.
\end{example}

\begin{example}[A rare large payoff]\label{ex:rare-payoff}
Let \(X=M\) with weight one and \(X=0\) with weight \(M-1\).  Then
\(\E[X]=1\) for every positive \(M\), even though the nonzero outcome has
probability \(1/M\).  This exact family anticipates the pointwise-convergence
counterexample in \Cref{sec:convergence}: a stable expectation does not by
itself provide concentration or typicality.
\end{example}

\begin{audit}[The article's die/Chebyshev percentage]\label{audit:die-cheb}
For a fair die, \(\mu=7/2\) and \(\sigma^2=35/12\).  Chebyshev with radius
\(5/2\) gives
\[
 \Pp\bigl(\abs{X-7/2}<5/2\bigr)
 \ge 1-\frac{35/12}{(5/2)^2}
 =1-\frac7{15}=\frac8{15}\approx53.3\%.
\]
The event \(\abs{X-7/2}<5/2\) is \(1<X<6\), whose exact probability is
\(4/6=2/3\).  By contrast, the inclusive event \(1\le X\le6\) has probability
one.  Thus the ``about 53 percent'' bound is mathematically valid for the open
interval but cannot describe the inclusive support event.
\end{audit}

\section{Algebraic properties, tails, and variance}\label{sec:properties}

\subsection{Linearity and order}

\begin{theorem}[Nonnegativity]\label{thm:nonneg}
If \(X_i\ge0\) at every positive-weight outcome, then \(\E[X]\ge0\).
\end{theorem}

\begin{proof}
Every weighted summand \(w_iX_i\) is nonnegative.  Their finite sum is
nonnegative, and the denominator \(W\) is positive.  With rational codes,
clear the common positive denominator of the \(X_i\) and prove the resulting
natural inequality by repeated addition.
\end{proof}

\begin{theorem}[Additivity and homogeneity]\label{thm:linearity}
For rational observables \(X,Y\) and \(a\in\Q\),
\[
 \E[X+Y]=\E[X]+\E[Y],\qquad
 \E[aX]=a\E[X].
\]
\end{theorem}

\begin{proof}
For additivity,
\[
 W\E[X+Y]=\sum_iw_i(X_i+Y_i)
 =\sum_iw_iX_i+\sum_iw_iY_i
 =W(\E[X]+\E[Y]).
\]
The middle equality is the bounded-sum addition theorem proved by
quantifier-free induction.  Since \(W>0\), the rational codes on both sides are
equal.  Homogeneity follows by pulling the fixed factor \(a\) through each
finite addition.  Internally one first chooses a common denominator for
\(a,X_i,Y_i\), so only signed natural polynomial identities remain.
\end{proof}

\begin{corollary}[Finite linear combinations]\label{cor:finite-linearity}
For a coded list \((a_j,X_j)_{j<r}\),
\[
 \E\left[\sum_{j<r}a_jX_j\right]
 =\sum_{j<r}a_j\E[X_j].
\]
\end{corollary}

\begin{proof}
Use quantifier-free induction on \(r\).  The induction formula is an equality
of coded rationals.  The successor step is \Cref{thm:linearity}.
\end{proof}

\begin{corollary}[Expectation as a linear form]\label{cor:linear-form}
For a fixed finite law, \(X\mapsto\E[X]\) is a rational linear functional on
the finite rational module of observables.
\end{corollary}

\begin{theorem}[Monotonicity]\label{thm:monotone}
If \(X_i\le Y_i\) at every positive-weight outcome, then
\(\E[X]\le\E[Y]\).
\end{theorem}

\begin{proof}
The difference \(Y-X\) is pointwise nonnegative, so
\(0\le\E[Y-X]=\E[Y]-\E[X]\) by nonnegativity and linearity.
Equivalently, sum the pointwise inequalities directly.  The second proof is
slightly shorter in the object language because it avoids introducing a
signed subtraction term.
\end{proof}

\begin{corollary}[Constants and iterated expectation]\label{cor:constant}
For constant \(c\), \(\E[c]=c\).  Consequently the expectation of the constant
random variable \(\E[X]\) is again \(\E[X]\).
\end{corollary}

\begin{proof}
\(\sum_iw_ic=c\sum_iw_i=cW\).  Apply the same equation with
\(c=\E[X]\).
\end{proof}

\subsection{Positive and negative parts}

\begin{definition}[Parts and absolute value]\label{def:parts}
For a coded rational \(x\), define
\[
 x^+=\max(x,0),\qquad x^-=\max(-x,0),\qquad \abs{x}=x^++x^-.
\]
These operations are primitive recursive because rational order is decidable.
Apply them pointwise to observables.
\end{definition}

\begin{lemma}[Pointwise part identities]\label{lem:part-identities}
For every coded rational \(x\),
\[
 x=x^+-x^-,\qquad \abs{x}=x^++x^-,\qquad x^+x^-=0.
\]
\end{lemma}

\begin{proof}
Decide whether \(x\ge0\).  In the first branch, \(x^+=x\) and \(x^-=0\); in
the second, \(x^+=0\) and \(x^-=-x\).  This is a permitted intuitionistic
case split because the coded comparison is decidable.
\end{proof}

\begin{theorem}[Expectation and parts]\label{thm:expect-parts}
For every finite rational observable,
\[
 \E[X]=\E[X^+]-\E[X^-],\qquad
 \E[\abs X]=\E[X^+]+\E[X^-].
\]
Moreover,
\[
 \abs{\E[X]}\le\E[\abs X].
\]
\end{theorem}

\begin{proof}
The two equalities follow from \Cref{lem:part-identities} and linearity.  Since
\(-\abs X\le X\le\abs X\), monotonicity gives
\(-\E\abs X\le\E X\le\E\abs X\), which is equivalent to the final
inequality by decidable rational order.
\end{proof}

\begin{theorem}[Zero absolute expectation]\label{thm:zero-abs}
If \(\E[\abs X]=0\), then \(X=0\) almost surely.  More generally, if
\(\E[\abs{X-Y}]=0\), then \(X=Y\) almost surely.
\end{theorem}

\begin{proof}
After choosing a common denominator \(D\), the hypothesis is
\[
 \sum_{i<m}w_i\abs{a_i}=0,
\]
where \(X_i=a_i/D\).  Each term is a natural number.  Zero-sum elimination
proves \(w_i\abs{a_i}=0\) for every \(i\).  Thus every positive-weight outcome
has \(a_i=0\).  Equivalently, the weighted failure indicator has sum zero.
Apply the same argument to \(X-Y\).
\end{proof}

\begin{theorem}[Almost-sure congruence]\label{thm:ae-congruence}
If \(X=Y\) almost surely, then \(\E[X]=\E[Y]\).  If
\(g:\Q\to\Q\) is a fixed primitive-recursive map on the finite input codes,
then also \(g(X)=g(Y)\) almost surely and \(\E[g(X)]=\E[g(Y)]\).
\end{theorem}

\begin{proof}
At every outcome with positive weight the values agree, by zero-sum
elimination on the exceptional mass.  Hence the weighted summands agree and so
do their bounded sums.  Equality is preserved by applying the function code
\(g\), giving the second assertion.
\end{proof}

\subsection{Indicators and event algebra}

\begin{theorem}[Indicator expectation]\label{thm:indicator}
For every decidable finite event \(A\),
\[
 \E[\ind_A]=\Pp(A).
\]
For events \(A,B\),
\[
 \ind_{A\cup B}=\ind_A+\ind_B-\ind_{A\cap B},
\]
and therefore
\[
 \Pp(A\cup B)=\Pp(A)+\Pp(B)-\Pp(A\cap B).
\]
If \(A\cap B=\varnothing\), the intersection term vanishes.
\end{theorem}

\begin{proof}
The first formula is the definition.  For the second, decide the four possible
truth-value pairs of \(A(i),B(i)\) at each finite outcome.  The displayed
identity holds in each case.  Apply expectation and linearity.  The formal
version may avoid signed subtraction by moving the intersection term to the
left.
\end{proof}

\begin{corollary}[Finite union bound]\label{cor:union-bound}
For a coded list of events \(A_0,\ldots,A_{r-1}\),
\[
 \Pp\left(\bigcup_{j<r}A_j\right)\le\sum_{j<r}\Pp(A_j).
\]
\end{corollary}

\begin{proof}
Pointwise,
\(\ind_{\cup_jA_j}\le\sum_j\ind_{A_j}\).  Sum and use monotonicity and finite
linearity.  The union indicator is computed by bounded Boolean recursion.
\end{proof}

\subsection{Integer tail formulas}

\begin{theorem}[Nonnegative integer tail sum]\label{thm:tail-sum}
Let \(X_i\in\N\) and suppose \(X_i\le B\) for all outcomes.  Then
\[
 X_i=\sum_{k<B}\ind_{X_i>k}
\]
pointwise, and hence
\[
 \E[X]=\sum_{k<B}\Pp(X>k).
\]
Equivalently,
\(\E[X]=\sum_{k=1}^{B}\Pp(X\ge k)\).
\end{theorem}

\begin{proof}
For a fixed natural \(x\le B\), the indicators
\(\ind_{x>0},\ldots,\ind_{x>B-1}\) consist of exactly \(x\) ones followed by
zeros.  Prove
\(x=\sum_{k<B}\ind_{x>k}\) by quantifier-free induction on \(x\), or by
induction on \(B\) with the invariant truncated at \(\min(x,B)\).  Substitute
\(x=X_i\), multiply by \(w_i\), sum over \(i\), and exchange the two finite
sums.  The alternative form is the index shift \(k\mapsto k-1\).
\end{proof}

\begin{theorem}[Signed integer CDF formula]\label{thm:signed-tail}
Suppose \(-A\le X_i\le B\) for natural \(A,B\).  Then
\[
 \E[X]
 =\sum_{k=0}^{B-1}\Pp(X>k)
  -\sum_{k=0}^{A-1}\Pp(X\le -k-1).
\]
\end{theorem}

\begin{proof}
Apply \Cref{thm:tail-sum} to \(X^+\) and \(X^-\), and use
\(X=X^+-X^-\).  For natural \(k\),
\(X^->k\) is equivalent to \(-X>k\), hence to \(X\le-k-1\) because \(X\) is
integer-valued.  All comparisons are decidable.
\end{proof}

\begin{theorem}[Finite layer cake]\label{thm:finite-layer-cake}
If \(X_i\ge0\) are rational with common denominator \(D\) and
\(X_i\le B/D\), then
\[
 \E[X]=\frac1D\sum_{k<B}\Pp(DX>k).
\]
This is the exact finite precursor of
\(\E[X]=\int_0^\infty\Pp(X>t)\,dt\).
\end{theorem}

\begin{proof}
Apply \Cref{thm:tail-sum} to the natural-valued observable \(DX\), then divide
the resulting identity by \(D>0\).
\end{proof}

\subsection{Variance and covariance}

\begin{definition}[Variance and covariance]\label{def:variance}
Let \(\mu=\E[X]\) and \(\nu=\E[Y]\).  Define
\[
 \Var(X)=\E[(X-\mu)^2],\qquad
 \Cov(X,Y)=\E[(X-\mu)(Y-\nu)].
\]
\end{definition}

\begin{theorem}[Second-moment identities]\label{thm:variance-identity}
For finite rational observables,
\[
 \Var(X)=\E[X^2]-\E[X]^2,
\]
\[
 \Cov(X,Y)=\E[XY]-\E[X]\E[Y],
\]
and
\[
 \Var(X+Y)=\Var(X)+\Var(Y)+2\Cov(X,Y).
\]
\end{theorem}

\begin{proof}
Expand the relevant squares and products and use finite linearity and the
constant rule.  For example,
\begin{align*}
\E[(X-\mu)^2]
 &=\E[X^2-2\mu X+\mu^2]\\
 &=\E[X^2]-2\mu\E[X]+\mu^2
 =\E[X^2]-\mu^2.
\end{align*}
The formal proof is a polynomial identity after replacing each expectation by
its weighted numerator.  ProveIt's ring-normalization layer is designed for
exactly such denominator-cleared identities.
\end{proof}

\begin{theorem}[Variance is nonnegative and detects constants]\label{thm:variance-zero}
\(\Var(X)\ge0\).  Moreover \(\Var(X)=0\) if and only if \(X=\E[X]\) almost
surely.
\end{theorem}

\begin{proof}
The square \((X-\mu)^2\) is pointwise nonnegative.  If its expectation is zero,
zero-sum elimination shows the square is zero at every positive-weight
outcome, hence \(X_i=\mu\).  The converse is immediate.
\end{proof}

\begin{theorem}[The mean minimizes squared error]\label{thm:mean-minimizes}
For every rational constant \(a\),
\[
 \E[(X-a)^2]=\Var(X)+(\E[X]-a)^2.
\]
Consequently \(a=\E[X]\) minimizes the expected squared error; equality at
another \(a\) implies \(a=\E[X]\).
\end{theorem}

\begin{proof}
Write \(X-a=(X-\mu)+(\mu-a)\).  After squaring, the cross term has expectation
\(2(\mu-a)\E[X-\mu]=0\).  The remaining terms give the identity.  The final
claim follows from nonnegativity of the square \((\mu-a)^2\).
\end{proof}

\section{Inequalities: Markov, Chebyshev, Jensen, and norm bounds}\label{sec:inequalities}

\subsection{Markov and Chebyshev}

\begin{theorem}[Markov inequality]\label{thm:markov}
Let \(X\ge0\) pointwise and let \(a>0\) be rational.  Then
\[
 \Pp(X\ge a)\le\frac{\E[X]}{a}.
\]
Equivalently, \(a\Pp(X\ge a)\le\E[X]\).
\end{theorem}

\begin{proof}
At each outcome,
\[
 a\ind_{X\ge a}\le X.
\]
Indeed, outside the event the left side is zero, and on the event it is \(a\),
which is at most \(X\).  Take expectations, use monotonicity and the indicator
formula, and divide by the positive \(a\).  The object-language target is the
cross-multiplied inequality, so no division function is required.
\end{proof}

\begin{corollary}[Power Markov]\label{cor:power-markov}
For natural \(r>0\), rational \(a>0\), and any \(X\),
\[
 \Pp(\abs X\ge a)\le\frac{\E[\abs X^r]}{a^r}.
\]
\end{corollary}

\begin{proof}
Apply \Cref{thm:markov} to the nonnegative variable \(\abs X^r\).  The events
\(\abs X^r\ge a^r\) and \(\abs X\ge a\) coincide because exponentiation by a
positive natural is order preserving on nonnegative rationals.  This last fact
is proved by induction on \(r\).
\end{proof}

\begin{theorem}[Chebyshev inequality]\label{thm:chebyshev}
For \(a>0\),
\[
 \Pp(\abs{X-\E[X]}\ge a)\le\frac{\Var(X)}{a^2}.
\]
\end{theorem}

\begin{proof}
Apply power Markov with \(r=2\) to \(X-\E[X]\).  Its second absolute moment is
exactly the variance.
\end{proof}

\begin{corollary}[At least three quarters within two standard deviations]
\label{cor:three-quarters}
If a positive rational \(s\) satisfies \(s^2=\Var(X)\), then
\[
 \Pp(\abs{X-\E[X]}<2s)\ge\frac34.
\]
The same denominator-free statement is
\(4\Var(X)\Pp(\abs{X-\E[X]}\ge2s)\le\Var(X)\), with the zero-variance case
handled separately.
\end{corollary}

\begin{proof}
Chebyshev gives probability at most \(1/4\) to the complementary event.
Subtract from one using the finite event identity.  A square root need not be
constructed internally if the theorem is stated with an input \(s\) and the
certificate \(s^2=\Var(X)\).
\end{proof}

\subsection{Constructive Jensen by a supporting-line certificate}

A bare classical definition of convexity quantifies over real coefficients.
For \(\ipra\), the most useful finite formulation is a checkable supporting
line at the particular mean under consideration.

\begin{definition}[Jensen certificate]\label{def:jensen-cert}
Let \(\phi\) be a primitive-recursive rational evaluator on the finite list of
values \(X_i\) and on \(\mu=\E[X]\).  A Jensen certificate at \(\mu\) is a
rational \(s\) such that for every \(i<m\),
\[
 \phi(X_i)\ge \phi(\mu)+s(X_i-\mu).
\]
\end{definition}

\begin{theorem}[Certified finite Jensen inequality]\label{thm:jensen}
If \(s\) is a Jensen certificate for \(\phi\) at \(\mu=\E[X]\), then
\[
 \phi(\E[X])\le\E[\phi(X)].
\]
\end{theorem}

\begin{proof}
Multiply the certificate inequality by \(w_i\), sum, and divide by \(W\):
\[
 \E[\phi(X)]
 \ge \phi(\mu)+s\E[X-\mu]
 =\phi(\mu).
\]
The vanishing of the final expectation is linearity and the constant rule.
All certificate checks are finite rational comparisons.
\end{proof}

\begin{corollary}[Square and absolute-value Jensen]\label{cor:jensen-basic}
For every finite rational \(X\),
\[
 \E[X]^2\le\E[X^2],\qquad
 \abs{\E[X]}\le\E[\abs X].
\]
\end{corollary}

\begin{proof}
For \(\phi(x)=x^2\), choose \(s=2\mu\); the certificate difference is
\[
 x^2-\mu^2-2\mu(x-\mu)=(x-\mu)^2\ge0.
\]
For \(\phi(x)=\abs x\), decide the sign of \(\mu\).  If \(\mu\ge0\), choose
\(s=1\) and use \(\abs x\ge x\); if \(\mu<0\), choose \(s=-1\) and use
\(\abs x\ge-x\).
\end{proof}

\begin{proposition}[Piecewise-linear rational Jensen]\label{prop:polygonal-jensen}
Suppose \(\phi\) is given by a finite rational convex polygon: rational
breakpoints \(b_0<\cdots<b_r\) and nondecreasing rational slopes
\(s_0\le\cdots\le s_{r-1}\).  Then a Jensen certificate at every rational
\(\mu\) in the domain is primitive-recursively computable, and
\Cref{thm:jensen} applies.
\end{proposition}

\begin{proof}
Bounded search locates the interval containing \(\mu\).  Choose its slope, or
at a breakpoint either adjacent slope.  On each side of \(\mu\), telescope the
polygonal increments.  Nondecreasing slopes show that the graph lies above the
chosen affine line.  The argument is a finite sum of products of nonnegative
rational segment lengths and slope differences.
\end{proof}

\begin{remark}[Smooth convex functions]
For a smooth function represented by regular rational approximants, the same
proof works with an approximate supporting slope and a certified error term.
At precision \(p\), use a rational slope whose supporting inequality is within
\(2^{-p-2}\), apply finite Jensen, and absorb the endpoint errors.  Thus the
usual Jensen theorem for effectively convex functions is an effective analytic
schema, not a new logical principle.
\end{remark}

\subsection{Cauchy--Schwarz as a finite sum of squares}

\begin{lemma}[Weighted Lagrange identity]\label{lem:lagrange}
For rational lists \(u_i,v_i\) and natural weights \(w_i\), let
\[
 A=\sum_iw_iu_i^2,\quad B=\sum_iw_iv_i^2,\quad C=\sum_iw_iu_iv_i.
\]
Then
\[
 AB-C^2
 =\sum_{0\le i<j<m}w_iw_j(u_iv_j-u_jv_i)^2.
\]
\end{lemma}

\begin{proof}
Expand the left side as a double sum:
\[
 AB-C^2=\sum_{i,j}w_iw_j(u_i^2v_j^2-u_iv_iu_jv_j).
\]
The diagonal terms vanish.  Pair the terms \((i,j)\) and \((j,i)\) for
\(i<j\); their sum is
\[
 w_iw_j(u_i^2v_j^2+u_j^2v_i^2-2u_iv_iu_jv_j)
 =w_iw_j(u_iv_j-u_jv_i)^2.
\]
Finite pairing and Fubini are justified by \Cref{cor:finite-fubini}.  The
identity is polynomial and can also be discharged by the ProveIt ring tool
after the double-sum lemmas have been instantiated.
\end{proof}

\begin{theorem}[Cauchy--Schwarz]\label{thm:cauchy-schwarz}
For finite rational observables,
\[
 \E[XY]^2\le\E[X^2]\E[Y^2].
\]
Consequently
\[
 \abs{\Cov(X,Y)}^2\le\Var(X)\Var(Y).
\]
\end{theorem}

\begin{proof}
Apply \Cref{lem:lagrange} with \(u_i=X_i\) and \(v_i=Y_i\).  The right side is
nonnegative, so \(C^2\le AB\).  Divide by \(W^2>0\).  For covariance, apply
the first inequality to \(X-\E[X]\) and \(Y-\E[Y]\).
\end{proof}

\begin{corollary}[Equality criterion]\label{cor:cs-equality}
Equality holds in Cauchy--Schwarz precisely when
\[
 w_iw_j(X_iY_j-X_jY_i)^2=0
\]
for every \(i<j\).  In particular, on the positive-weight support the two
vectors are rationally proportional unless both vanish in the relevant
coordinates.
\end{corollary}

\begin{proof}
The difference in \Cref{lem:lagrange} is a finite sum of nonnegative rational
terms.  It is zero exactly when every term is zero, by zero-sum elimination.
The proportionality reading follows by selecting any positive-weight index
where one coordinate is nonzero and cross-multiplying.
\end{proof}

\subsection{Holder, Lyapunov, and Minkowski}

The article lists several norm inequalities.  For integer exponents one can
prove them in exact algebraic form; for arbitrary computable real exponents we
state the regular-real version with root and power moduli.

\begin{theorem}[Finite Holder with effective rational exponents]
\label{thm:holder}
Let \(p>1\) and \(q>1\) be fixed positive rational exponents satisfying
\(1/p+1/q=1\).  Assume the presentation includes regular evaluators for the
rational powers and roots, together with a checked Young-inequality certificate
at the requested precision.  Then
\[
 \E[\abs{XY}]
 \le \E[\abs X^p]^{1/p}\E[\abs Y^q]^{1/q}.
\]
The special case \(p=q=2\) is the exact rational theorem
\Cref{thm:cauchy-schwarz}.  For general rational conjugates, the statement is
an effective regular-real schema.  Its finite object-language core says that
rational upper bounds \(A,B>0\) satisfying
\(A^p\ge\E\abs X^p\) and \(B^q\ge\E\abs Y^q\), to certified precision, imply
\(\E\abs{XY}\le AB\) up to the allocated dyadic error.
\end{theorem}

\begin{proof}
Normalize by positive rational upper bounds \(A,B\).  Young's inequality
\(uv\le u^p/p+v^q/q\) for nonnegative inputs is verified at finite precision by
the supplied power certificates; for fixed integer powers it can instead be
reduced to a polynomial factorization after clearing denominators.  Apply it to
\(u=\abs X/A\) and \(v=\abs Y/B\), multiply by \(AB\), and take expectations.
The normalized moments are at most one, yielding the claimed bound, with the
certificate errors absorbed by guard bits.  Regular roots are found by
bisection, so their interval widths provide a primitive-recursive output
modulus.
\end{proof}

\begin{corollary}[Lyapunov moment monotonicity]\label{cor:lyapunov}
For \(0<r<s\) given by fixed positive rational exponents and a finite rational
observable,
\[
 \E[\abs X^r]^{1/r}\le\E[\abs X^s]^{1/s}.
\]
\end{corollary}

\begin{proof}
Apply Holder to \(\abs X^r\cdot1\) with conjugate exponents
\(s/r\) and \(s/(s-r)\), then use \(\E[1]=1\).  In the denominator-cleared
form, all roots appear only through certified rational enclosures.
\end{proof}

\begin{theorem}[Finite Minkowski]\label{thm:minkowski}
For fixed computable \(p\ge1\),
\[
 \E[\abs{X+Y}^p]^{1/p}
 \le \E[\abs X^p]^{1/p}+\E[\abs Y^p]^{1/p}.
\]
For \(p=1\) the theorem is exact rational algebra; for \(p=2\) it follows
from Cauchy--Schwarz; for general fixed rational \(p\) it is an effective
regular-real schema based on Holder.
\end{theorem}

\begin{proof}
For \(p=1\), use the pointwise triangle inequality and monotonicity.  For
\(p>1\), put \(q=p/(p-1)\).  The pointwise inequality
\[
 \abs{X+Y}^p
 \le \abs X\abs{X+Y}^{p-1}+\abs Y\abs{X+Y}^{p-1}
\]
followed by expectation and Holder gives
\[
 M^p\le (A+B)M^{p-1},
\]
where \(M=\norm{X+Y}_p\), \(A=\norm X_p\), and \(B=\norm Y_p\).
If \(M=0\) the claim is immediate; otherwise cancel the positive
\(M^{p-1}\).  Constructively, cancellation is performed on rational interval
bounds separated from zero; when separation is unavailable, the target error
already dominates \(M\), so the approximate inequality closes.  This standard
case split is decidable at each finite precision.
\end{proof}

\section{Independence, products, and conditional expectation}\label{sec:independence}

\subsection{Product laws}

\begin{definition}[Finite product law]\label{def:product-law}
Let \(L=(m,w,W)\) and \(M=(n,v,V)\).  Their product law has outcomes
\((i,j)\in\Fin(m)\times\Fin(n)\), weight \(w_iv_j\), and total weight
\[
 \sum_{i<m}\sum_{j<n}w_iv_j=WV.
\]
An observable depending only on the first coordinate is written \(X(i)\), and
one depending only on the second is written \(Y(j)\).
\end{definition}

\begin{theorem}[Multiplicativity under independence]\label{thm:ind-product}
On the product law,
\[
 \E[X(i)Y(j)]=\E_L[X]\E_M[Y].
\]
More generally, for finite lists of independent coordinates and rational
functions \(g_r\) of one coordinate,
\[
 \E\left[\prod_{r<k}g_r(X_r)\right]
 =\prod_{r<k}\E[g_r(X_r)].
\]
\end{theorem}

\begin{proof}
For two coordinates,
\begin{align*}
 WV\E[XY]
 &=\sum_{i<m}\sum_{j<n}w_iv_jX_iY_j\\
 &=\left(\sum_{i<m}w_iX_i\right)
   \left(\sum_{j<n}v_jY_j\right),
\end{align*}
by finite distributivity and Fubini.  Divide by \(WV>0\).  The multi-coordinate
formula follows by quantifier-free induction on \(k\), using a mixed-radix code
for product outcomes and the two-coordinate identity at each successor step.
\end{proof}

\begin{corollary}[Variance of an independent sum]\label{cor:ind-var-sum}
If \(X,Y\) are independent, then \(\Cov(X,Y)=0\) and
\[
 \Var(X+Y)=\Var(X)+\Var(Y).
\]
For a finite independent family,
\[
 \Var\left(\sum_{j<r}X_j\right)=\sum_{j<r}\Var(X_j).
\]
\end{corollary}

\begin{proof}
Multiplicativity gives \(\E[XY]=\E[X]\E[Y]\), so covariance is zero.  Insert
this into the variance expansion.  Induct on the number of summands; a new
coordinate is independent of every function of the preceding product
coordinates, a fact verified by another finite product factorization.
\end{proof}

\begin{example}[Expectation is not multiplicative in general]\label{ex:not-multiplicative}
Let \(X=Y\) be uniform on \(\{-1,1\}\).  Then
\(\E[X]=\E[Y]=0\), while \(XY=1\) identically and \(\E[XY]=1\).
Thus multiplicativity is a consequence of independence, not a basic property
of expectation.
\end{example}

\begin{example}[Uncorrelated does not imply independent]\label{ex:uncorrelated}
Let \(X\) be uniform on \(\{-1,0,1\}\) and let \(Y=X^2\).  Then
\[
 \E[X]=0,
 \qquad \E[XY]=\E[X^3]=0,
 \qquad \Cov(X,Y)=0.
\]
But \(Y\) is a deterministic function of \(X\), and the event \(Y=0\) is
exactly the event \(X=0\); the joint law does not factor.  All values are
finite rationals, so this counterexample is internal to \(\ipra\).
\end{example}

\subsection{Finite partitions and conditional expectation}

\begin{definition}[Finite conditioning partition]\label{def:partition}
A finite partition is a primitive-recursive map
\(c:\Fin(m)\to\Fin(r)\).  Its cell masses and weighted \(X\)-sums are
\[
 M_j=\sum_{i<m}w_i\ind_{c(i)=j},\qquad
 S_j=\sum_{i<m}w_iX_i\ind_{c(i)=j}.
\]
A conditional-expectation code is a cellwise rational value
\[
 h_j=\begin{cases}S_j/M_j,&M_j>0,\\0,&M_j=0.\end{cases}
\]
The observable \(\E[X\mid c]\) takes value \(h_{c(i)}\) at outcome \(i\).
\end{definition}

The arbitrary value on a zero-mass cell is harmless: its total weight is zero.
The branch \(M_j=0\) is decidable, and division is only performed in the
positive branch.

\begin{theorem}[Cell integral characterization]\label{thm:cell-characterization}
For every cell \(j<r\),
\[
 \sum_{i<m}w_i\E[X\mid c](i)\ind_{c(i)=j}
 =\sum_{i<m}w_iX_i\ind_{c(i)=j}.
\]
Equivalently,
\[
 \E[\E[X\mid c]\ind_{c=j}]=\E[X\ind_{c=j}].
\]
\end{theorem}

\begin{proof}
If \(M_j=0\), zero-sum elimination shows every weight in the cell is zero, so
both sides vanish.  If \(M_j>0\), the left numerator is
\(h_jM_j=(S_j/M_j)M_j=S_j\).  The proof is a decidable split on the natural
number \(M_j\).
\end{proof}

\begin{theorem}[Law of total expectation]\label{thm:total-expectation}
For every finite partition,
\[
 \E[\E[X\mid c]]=\E[X].
\]
More generally, if \(g\) is constant on each cell, then
\[
 \E[g\,\E[X\mid c]]=\E[gX].
\]
\end{theorem}

\begin{proof}
Sum the cell characterization over \(j<r\).  For each outcome exactly one cell
indicator is one, so both double sums collapse to the corresponding full sum.
For the generalized identity, multiply the equation for cell \(j\) by the
cell value \(g_j\) before summing.
\end{proof}

\begin{theorem}[Tower property for nested finite partitions]\label{thm:tower}
Suppose \(d:\Fin(m)\to\Fin(s)\) refines
\(c:\Fin(m)\to\Fin(r)\): there is \(q:\Fin(s)\to\Fin(r)\) with
\(c=q\circ d\).  Then
\[
 \E[\E[X\mid d]\mid c]=\E[X\mid c]
\]
almost surely.
\end{theorem}

\begin{proof}
Both sides are constant on each \(c\)-cell.  By
\Cref{thm:cell-characterization,thm:total-expectation}, their weighted integral
on that cell equals the integral of \(X\) on the cell.  If the cell has
positive mass, equality of integral and constancy imply equality of the two
cell values after cancellation.  If the cell has zero mass, disagreement is
confined to a null cell.  The refinement relation makes the regrouping of
\(d\)-cells into a \(c\)-cell a finite Fubini identity.
\end{proof}

\begin{theorem}[Conditional Jensen with certificates]\label{thm:conditional-jensen}
Let \(\phi\) admit a finite Jensen certificate at the conditional mean in each
positive-mass cell.  Then
\[
 \phi(\E[X\mid c])\le \E[\phi(X)\mid c]
\]
almost surely.
\end{theorem}

\begin{proof}
Apply certified Jensen separately to the normalized finite law on each
positive-mass cell.  A zero-mass cell may be ignored.  Combining the finite
list of cellwise inequalities yields the almost-sure statement.
\end{proof}

\subsection{Fubini and effective products}

Finite Fubini is exact.  For countable or continuous product presentations, a
Fubini theorem needs an absolute-integrability modulus.  The constructive
pattern is as follows: truncate both coordinates at stages computed from the
tail modulus; apply finite Fubini; and bound the discarded absolute mass.
This will be formalized in \Cref{sec:integration}.  Independence then implies
multiplicativity whenever the product and the two factors are effectively
integrable.

\section{Samples, averages, and finitary laws of large numbers}\label{sec:samples}

\subsection{Sample means}

Let \(X_0,\ldots,X_{n-1}\) be coordinates on a finite product law.  Their
sample mean is
\[
 \overline X_n=\frac1n\sum_{j<n}X_j
\]
for \(n>0\).

\begin{theorem}[Unbiased sample mean]\label{thm:sample-mean}
If each \(X_j\) has expectation \(\mu\), then
\[
 \E[\overline X_n]=\mu.
\]
No independence assumption is needed.
\end{theorem}

\begin{proof}
By finite linearity,
\[
 \E[\overline X_n]
 =\frac1n\sum_{j<n}\E[X_j]
 =\frac1n(n\mu)=\mu.
\]
In \(\ipra\), the identity \(\sum_{j<n}\mu=n\mu\) is proved by
quantifier-free induction on \(n\), and the theorem is kept in
cross-multiplied form \(n\E[\overline X_n]=n\mu\).
\end{proof}

\begin{theorem}[Variance of the independent sample mean]\label{thm:sample-var}
If the coordinates are independent and identically distributed with variance
\(\sigma^2\), then
\[
 \Var(\overline X_n)=\frac{\sigma^2}{n}.
\]
\end{theorem}

\begin{proof}
By independent variance additivity,
\(\Var(\sum_{j<n}X_j)=n\sigma^2\).  Homogeneity of variance,
\(\Var(aX)=a^2\Var(X)\), follows from its definition.  Therefore
\[
 \Var(\overline X_n)=\frac1{n^2}n\sigma^2=\frac{\sigma^2}{n}.
\]
All calculations are rational when the underlying finite law is rational.
\end{proof}

\begin{theorem}[Expected residual sum of squares]\label{thm:residual}
For independent identically distributed coordinates with variance
\(\sigma^2\),
\[
 \E\left[\sum_{j<n}(X_j-\overline X_n)^2\right]=(n-1)\sigma^2.
\]
Consequently, for \(n>1\),
\(\frac1{n-1}\sum_{j<n}(X_j-\overline X_n)^2\) is an unbiased estimator of
\(\sigma^2\).
\end{theorem}

\begin{proof}
The pointwise polynomial identity
\[
 \sum_{j<n}(X_j-\overline X_n)^2
 =\sum_{j<n}(X_j-\mu)^2-n(\overline X_n-\mu)^2
\]
is proved by expanding both sides and using
\(\sum_j(X_j-\overline X_n)=0\).  Take expectations.  The first term has
expectation \(n\sigma^2\); the second has expectation
\(n\Var(\overline X_n)=\sigma^2\).  Subtract.
\end{proof}

\subsection{Weak law as a finite quantitative theorem}

\begin{theorem}[Chebyshev weak-law bound]\label{thm:weak-law}
Under the assumptions of \Cref{thm:sample-var}, for rational
\(\varepsilon>0\),
\[
 \Pp(\abs{\overline X_n-\mu}\ge\varepsilon)
 \le\frac{\sigma^2}{n\varepsilon^2}.
\]
Given rational \(\delta>0\), any natural \(n\) with
\(n\varepsilon^2\delta\ge\sigma^2\) guarantees probability at most
\(\delta\).
\end{theorem}

\begin{proof}
Apply Chebyshev to \(\overline X_n\), insert
\(\Var(\overline X_n)=\sigma^2/n\), and clear the positive denominators.  The
last sentence is merely transitivity of rational order.  It is a direct
\(\ipra\) contract: the input supplies \(n,\varepsilon,\delta,\sigma^2\) and
a checker verifies one natural inequality.
\end{proof}

\subsection{Kolmogorov's finite maximal inequality}

For a finite independent centered sequence put
\(S_k=\sum_{j<k}X_j\).  Every product sample space is finite, so first-crossing
events are decidable by bounded search.

\begin{theorem}[Finite Kolmogorov maximal inequality]\label{thm:kolmogorov}
Let \(X_0,\ldots,X_{n-1}\) be independent finite rational observables with
\(\E[X_j]=0\).  For rational \(a>0\),
\[
 \Pp\left(\max_{1\le k\le n}\abs{S_k}\ge a\right)
 \le\frac{\sum_{j<n}\Var(X_j)}{a^2}.
\]
\end{theorem}

\begin{proof}
Let \(A_k\) be the event that \(k\) is the first index with
\(\abs{S_k}\ge a\).  The \(A_k\) are disjoint and their union is the maximal
event.  On \(A_k\), \(S_k^2\ge a^2\), so
\[
 a^2\Pp(A_k)\le\E[S_k^2\ind_{A_k}].
\]
Write \(S_n=S_k+(S_n-S_k)\).  Since \(A_k\) depends only on the first \(k\)
coordinates and the tail increment is independent of them with mean zero,
\[
 \E[S_k(S_n-S_k)\ind_{A_k}]=0.
\]
Also \(\E[(S_n-S_k)^2\ind_{A_k}]\ge0\).  Hence
\[
 \E[S_n^2\ind_{A_k}]\ge\E[S_k^2\ind_{A_k}]
 \ge a^2\Pp(A_k).
\]
Sum over the disjoint first-crossing events:
\[
 a^2\Pp\left(\max_{k\le n}\abs{S_k}\ge a\right)
 \le\E\left[S_n^2\sum_k\ind_{A_k}\right]
 \le\E[S_n^2].
\]
The final second moment is \(\sum_{j<n}\Var(X_j)\) by independence and
centering.  Every event and every sum in this proof is finite.  The only
induction is the construction of partial sums and the bounded search for the
first crossing.
\end{proof}

\subsection{A finitary strong-law contract}

A qualitative statement ``\(\overline X_n\to\mu\) almost surely'' quantifies
over an infinite sequence and does not by itself expose a usable rate.  The
following theorem is the precise finite-error content needed by many
applications.

\begin{theorem}[Dyadic eventual-uniform bound]\label{thm:dyadic-slln}
Let \(X_0,X_1,\ldots\) be presented as an iid finite rational law with mean
\(\mu\) and variance \(\sigma^2\).  For naturals \(K\) and \(H\ge2^K\), and
rational \(\varepsilon>0\),
\[
 \Pp\left(\exists n\in[2^K,H]\;:
   \abs{\overline X_n-\mu}\ge 2\varepsilon\right)
 \le \frac{4\sigma^2}{\varepsilon^2 2^K}.
\]
Consequently, if
\(4\sigma^2\le\delta\varepsilon^2 2^K\), then with probability at least
\(1-\delta\), every \(n\in[2^K,H]\) satisfies
\(\abs{\overline X_n-\mu}<2\varepsilon\).
\end{theorem}

\begin{proof}
Center the variables and write \(S_n=\sum_{j<n}(X_j-\mu)\).  Partition
\([2^K,H]\) into dyadic blocks \([2^r,2^{r+1})\).  If for some
\(n\in[2^r,2^{r+1})\) we have \(\abs{S_n}\ge2\varepsilon n\), then certainly
\(\max_{k\le2^{r+1}}\abs{S_k}\ge2\varepsilon2^r\).  Kolmogorov's inequality
gives
\[
 \Pp(\text{bad block }r)
 \le\frac{2^{r+1}\sigma^2}{(2\varepsilon2^r)^2}
 =\frac{\sigma^2}{2\varepsilon^2 2^r}.
\]
The union bound over \(r\ge K\) up to the last block is at most the infinite
geometric majorant
\[
 \sum_{r=K}^{\infty}\frac{\sigma^2}{2\varepsilon^2 2^r}
 =\frac{\sigma^2}{\varepsilon^2 2^K},
\]
which is stronger than the displayed constant \(4\).  We retain \(4\) to
absorb endpoint conventions, non-power-of-two horizons, and a version obtained
by controlling block increments rather than maxima from zero.  Formally, for a
fixed finite \(H\), only a finite geometric sum is used, and it is bounded by
\(2^{1-K}\) through an elementary induction.

Finally, \(\abs{\overline X_n-\mu}=\abs{S_n}/n\).  The complement of the bad
event therefore gives the simultaneous bound.  The theorem is a family of
finite \(\ipra\) statements uniform in \(K,H,\varepsilon,\delta\).
\end{proof}

\begin{remark}[Relation to the classical strong law]
The family in \Cref{thm:dyadic-slln} has a tail bound independent of the finite
horizon \(H\).  Metamathematically, letting \(H\) grow and then \(K\) grow
recovers the usual almost-sure convergence interpretation.  Inside \(\ipra\),
the directly usable theorem is the finite contract: given an error tolerance,
failure budget, and finite observation horizon, compute and verify a burn-in
index.  This avoids silently importing countable additivity or an infinitary
``eventually'' predicate.
\end{remark}

\section{Countable laws and effective series}\label{sec:countable}

A countable probability law cannot be inserted into first-order arithmetic as
an actual completed set.  It is presented by primitive-recursive rational
terms and moduli.  The finite partial sums are exact objects of the preceding
sections.

\subsection{Regular sums from absolute-tail moduli}

\begin{definition}[Absolutely summable rational series name]\label{def:abs-series}
An absolute-series name consists of a primitive-recursive rational sequence
\(a_n\) and a primitive-recursive modulus \(N(p)\) such that
\[
 \sum_{n=N(p)}^{M}\abs{a_n}\le2^{-p}
 \qquad\text{for every }M\ge N(p).
\]
The value of the series is represented by the rational intervals
\[
 I_p=\left[\sum_{n<N(p+1)}a_n-2^{-(p+1)},
          \sum_{n<N(p+1)}a_n+2^{-(p+1)}\right].
\]
\end{definition}

The intervals are nested after a harmless monotone closure operation, and
their widths are at most \(2^{-p}\).  Thus they form a regular-real name.  All
interval endpoints and all checks are primitive recursive.

\begin{theorem}[Effective Cauchy convergence]\label{thm:abs-convergence}
Every absolute-series name determines a unique regular real.  Addition,
scalar multiplication, and finite regrouping of such names preserve effective
summability, with moduli computed from the input moduli.
\end{theorem}

\begin{proof}
If \(M,L\ge N(p)\), then
\[
 \abs{\sum_{n<M}a_n-\sum_{n<L}a_n}
 \le\sum_{n=\min(M,L)}^{\max(M,L)-1}\abs{a_n}
 \le2^{-p}.
\]
This proves the Cauchy estimate using only a finite triangle inequality.
For sums of two series, request one additional bit from each tail; for scalar
multiplication by \(c\), request enough extra bits that
\(\abs c\,2^{-q}\le2^{-p}\).  Finite regrouping is exact at every sufficiently
large finite stage.
\end{proof}

\begin{definition}[Countable probability and expectation]\label{def:countable-law}
A countable rational law is a primitive-recursive sequence
\(p_n\ge0\) with an absolute-tail modulus and regular sum one.  A rational
observable \(x_n\) is integrable when \(a_n=p_nx_n\) has an absolute-tail
modulus.  Its expectation is the regular sum
\[
 \E[X]=\sum_{n=0}^{\infty}p_nx_n.
\]
\end{definition}

\begin{theorem}[Effective invariance under reindexing]\label{thm:reindex}
Let \(a_n\) have an absolute-tail modulus.  Suppose
\(\pi:\N\to\N\) is a primitive-recursive bijection equipped with a coverage
modulus \(C(N)\) satisfying
\[
 \{0,\ldots,N-1\}\subseteq
 \{\pi(0),\ldots,\pi(C(N)-1)\}.
\]
Then \(\sum_na_{\pi(n)}\) has the same regular sum as \(\sum_na_n\), with a
primitive-recursively derived modulus.
\end{theorem}

\begin{proof}
Choose \(N=N(p+2)\) from the absolute-tail modulus and
\(K=C(N)\).  The first \(K\) permuted terms contain every original term below
\(N\).  Their difference from the first \(N\) original terms is a finite sum
of terms whose indices are at least \(N\), hence has absolute value at most
\(2^{-(p+2)}\).  Any additional permuted tail is bounded by the same absolute
tail.  The triangle inequality gives the requested \(2^{-p}\) enclosure.
No appeal to a completed permutation set is made; the finite coverage
certificate is the essential constructive data.
\end{proof}

\subsection{Conditional convergence and rearrangement}

\begin{proposition}[Alternating harmonic series]\label{prop:alt-harmonic}
The alternating harmonic partial sums
\[
 A_N=\sum_{n=1}^{N}\frac{(-1)^{n+1}}n
\]
form a regular real with error at most \(1/(N+1)\).  This regular real is the
effectively defined \(\log 2\).
\end{proposition}

\begin{proof}
Group adjacent terms:
\[
 \frac1{2k-1}-\frac1{2k}=\frac1{(2k-1)2k}>0.
\]
The even partial sums increase and the odd partial sums decrease.  Their gap is
the next term, at most \(1/(N+1)\).  Given precision \(p\), bounded search finds
an \(N\) with \(N+1\ge2^p\), producing a regular interval.  We may define
\(\log2\) by this name; equivalence with integral or Taylor definitions is a
separate effective analytic theorem.
\end{proof}

\begin{theorem}[A concrete rearrangement changes the sum]\label{thm:rearrange}
Rearrange the alternating harmonic series in blocks containing two positive
terms followed by one negative term:
\[
 1+\frac13-\frac12+\frac15+\frac17-\frac14+\cdots.
\]
The resulting block partial sums converge effectively to
\(\tfrac32\log2\), not \(\log2\).
\end{theorem}

\begin{proof}
After \(N\) complete blocks, the sum is
\[
 R_N=\sum_{k=1}^{2N}\frac1{2k-1}-\frac12H_N
 =H_{4N}-\frac12H_{2N}-\frac12H_N,
\]
where \(H_n=\sum_{k=1}^n1/k\).  Write
\(H_m=\log m+\gamma+\epsilon_m\), but do not assume an abstract constant
\(\gamma\): use the elementary finite bounds obtained by comparing
\(1/k\) with \(\int_k^{k+1}1/x\,dx\), which show that
\(H_m-\log m\) is a Cauchy name and that the three error terms cancel with an
explicit \(O(1/N)\) bound.  The logarithmic part is
\[
 \log(4N)-\tfrac12\log(2N)-\tfrac12\log N
 =\tfrac32\log2.
\]
Compact rational quadrature for \(1/x\) supplies the needed regular log name
and modulus.  Equivalently, one can compare \(R_N\) directly with three
alternating-harmonic subsequences.  Every comparison is finite at precision
\(p\), with \(N\) chosen larger than a primitive-recursive multiple of
\(2^p\).

The original series is not absolutely summable because its absolute partial
sums are \(H_N\), which exceed any prescribed natural bound by a primitive-
recursive block estimate.  Thus \Cref{thm:reindex} does not apply, exactly as it
should not.
\end{proof}

\subsection{The logarithmic probability example}

\begin{proposition}[A countable law normalized by \(\log2\)]
\label{prop:log-law}
For \(n\ge1\), put
\[
 p_n=\frac{1}{\log2}\frac1{n2^n}.
\]
Then \(\sum_{n\ge1}p_n=1\), and the integer-valued variable \(X(n)=n\) has
\[
 \E[X]=\frac1{\log2}.
\]
Both statements have explicit primitive-recursive tail moduli.
\end{proposition}

\begin{proof}
The finite identity
\[
 \sum_{n=1}^{N}\frac{x^n}{n}
\]
after differentiation is the geometric sum
\(\sum_{n=1}^{N}x^{n-1}=(1-x^N)/(1-x)\).  Integrating its rational polynomial
identity from \(0\) to \(1/2\), with exact polynomial antiderivatives, yields
\[
 \sum_{n=1}^{N}\frac1{n2^n}
 =\int_0^{1/2}\frac{1-t^N}{1-t}\,dt.
\]
The omitted tail is bounded by
\[
 \sum_{n>N}\frac1{n2^n}
 \le\frac1{N+1}\sum_{n>N}2^{-n}
 =\frac{2^{-N}}{N+1}.
\]
The integral limit is the regular \(\log2\).  Hence the \(p_n\) sum to one.
For the expectation,
\[
 \sum_{n=1}^{N}np_n
 =\frac1{\log2}\sum_{n=1}^{N}2^{-n},
\]
and the geometric tail is exactly \(2^{-N}\).  Taking the effective limit gives
\(1/\log2\); the reciprocal is computed by interval bisection after a rational
positive lower bound for \(\log2\) has been certified.
\end{proof}

\subsection{Extended nonnegative expectation}

\begin{definition}[Lower-real expectation]\label{def:lower-expect}
For a nonnegative countable observable, define the increasing rational partial
expectations
\[
 E_N=\sum_{n<N}p_nx_n.
\]
The sequence \((E_N)\) is a lower-real name.  It is finite as a regular real if
an upper-tail modulus is supplied; it is \(+\infty\) when a primitive-recursive
divergence function \(D(B)\) satisfies \(E_{D(B)}\ge B\) for every natural
\(B\).
\end{definition}

\begin{proposition}[St Petersburg game]\label{prop:st-petersburg}
Let the payoff be \(2^n\) with probability \(2^{-n}\) for \(n\ge1\).  The
partial expectation through outcome \(N\) is exactly \(N\); therefore the
extended expectation is \(+\infty\), with divergence function \(D(B)=B\).
\end{proposition}

\begin{proof}
Each term contributes \(2^{-n}2^n=1\).  Thus the first \(N\) terms sum to
\(N\).  This is a quantifier-free induction on \(N\).  The probabilities
normalize because \(\sum_{n=1}^{N}2^{-n}=1-2^{-N}\) with an explicit tail.
\end{proof}

\begin{definition}[Signed extended expectation]\label{def:signed-extended}
For a signed countable observable, form the lower-real expectations of
\(X^+\) and \(X^-\).  If both are finite, define
\(\E[X]=\E[X^+]-\E[X^-]\).  If exactly one is infinite, define the corresponding
signed infinity.  If both are infinite, the expectation is undefined.
\end{definition}

\begin{proposition}[Classification by positive and negative parts]
\label{prop:parts-classification}
The four cases in \Cref{def:signed-extended} are disjoint and exhaustive for a
presentation whose two lower-real names are each certified finite or divergent.
In the finite-finite case, the resulting regular real agrees with the absolute
series \(\sum p_nx_n\), and \(\sum p_n\abs{x_n}\) is finite.
\end{proposition}

\begin{proof}
The certificates make each part decidably classified at the level of the
presentation.  In the finite-finite case,
\(\abs{x}=x^++x^-\) termwise, so the absolute expectation is the sum of two
finite regular reals.  Conversely, an absolute-tail modulus bounds both parts.
The signed identity follows from finite partial sums and passage to regular
limits using their moduli.
\end{proof}

\begin{proposition}[Undefined alternating unbounded example]
\label{prop:undefined-alt}
Let \(p_n=c/n^2\) for \(n\ge1\), where \(c\) is the regular reciprocal of
\(\sum_{n\ge1}1/n^2\), and let \(X(n)=(-1)^{n+1}n\).  Then both
\(\E[X^+]\) and \(\E[X^-]\) diverge, so \(\E[X]\) is undefined even though
symmetric truncations may converge.
\end{proposition}

\begin{proof}
For odd \(n\), the positive contribution is \(c/n\); for even \(n\), the
negative contribution is \(c/n\).  The odd and even harmonic subseries both
diverge.  A primitive-recursive divergence witness follows from grouping
\(2^{r-1}<n\le2^r\): each full block contributes a fixed positive rational
lower bound.  The constant \(c>0\) has a certified rational lower bound from a
finite upper bound on \(\sum1/n^2\).  Thus each part exceeds any prescribed
bound after a computable number of blocks.

The Basel value \(\pi^2/6\) is unnecessary for divergence; only a positive
normalizing regular real is used.  A wholly rational variant takes
\(p_n=1/(n(n+1))\) and the same alternating magnitude \(n\), leading to the
same two divergent parts.
\end{proof}

\subsection{Countable tail sums}

\begin{corollary}[Effective integer tail formula]\label{cor:countable-tail}
Let \(X\ge0\) be integer-valued under a countable law.  Suppose the expectation
tail has a modulus.  Then
\[
 \E[X]=\sum_{k=0}^{\infty}\Pp(X>k)
\]
as equality of regular or lower reals.
\end{corollary}

\begin{proof}
Truncate both the outcome index and the value at finite bounds.  Apply the exact
finite tail formula.  Monotonicity makes the two truncation directions
increasing.  The supplied expectation-tail modulus controls the discarded
weighted values; this also controls the tail-probability double sum because it
is termwise identical at every finite stage.  Hence the regular names agree.
Without a finiteness modulus, the same argument gives equality of lower reals.
\end{proof}

\section{Effective integration and continuous laws}\label{sec:integration}

\subsection{Regular reals and simple-law names}

\begin{definition}[Regular rational interval name]\label{def:regular}
A regular real name is a pair of primitive-recursive rational functions
\(\ell(p),u(p)\) such that
\[
 \ell(p)\le u(p),\qquad
 [\ell(p+1),u(p+1)]\subseteq[\ell(p),u(p)],\qquad
 u(p)-\ell(p)\le2^{-p}.
\]
Arithmetic is performed by rational interval operations with extra guard bits.
\end{definition}

\begin{definition}[Effective simple random variable]\label{def:simple-name}
An effective \(L^1\) simple-law name is a primitive-recursive sequence of
finite rational laws \(X^{(p)}\), together with couplings of any two levels and
a modulus
\[
 \E\abs{X^{(p)}-X^{(q)}}\le2^{-p}
 \qquad(q\ge p).
\]
The couplings are finite common refinements: a finite weighted law carrying
both approximants and having the prescribed marginals.
\end{definition}

This presentation is deliberately operational.  It does not quantify over a
set of measurable functions.  Each level is a finite object, and the stated
\(L^1\) inequality is a rational certificate.

\begin{theorem}[Expectation of an \(L^1\) simple-law name]
\label{thm:l1-expect}
The rational sequence \(e_p=\E[X^{(p)}]\) is a regular Cauchy name.  More
precisely, for \(q\ge p\),
\[
 \abs{e_q-e_p}\le2^{-p}.
\]
Thus \(\E[X]\) is constructively defined.  Equivalent names at zero effective
\(L^1\) distance have the same expectation.
\end{theorem}

\begin{proof}
On the supplied common refinement,
\[
 \abs{e_q-e_p}
 =\abs{\E[X^{(q)}-X^{(p)}]}
 \le\E\abs{X^{(q)}-X^{(p)}}
 \le2^{-p}.
\]
The first inequality is the exact finite absolute-expectation inequality.
Intervals centered at \(e_{p+2}\) with radius \(2^{-(p+2)}\), followed by a
nested-hull normalization, give \Cref{def:regular}.  If two names have
arbitrarily small effective \(L^1\) distance, apply the same estimate across a
common refinement.
\end{proof}

\begin{theorem}[Lifted algebra]\label{thm:l1-algebra}
Finite linearity, monotonicity, almost-sure/zero-\(L^1\) congruence, the absolute
expectation bound, and certified Jensen lift from finite laws to effective
\(L^1\) names.  The output modulus is primitive-recursively computed from the
input moduli and coefficient bounds.
\end{theorem}

\begin{proof}
At output precision \(p\), choose a common finite level with enough guard bits.
Apply the corresponding exact finite theorem.  Bound the discrepancy between
that finite result and each limiting name by the \(L^1\) modulus and
\(\abs{\E Z}\le\E\abs Z\).  For a coefficient \(c\), request \(q\) such that
\(\abs c2^{-q}\le2^{-p}\).  Monotonicity is stated with an effective
nonnegative-difference presentation, rather than an uncheckable pointwise
order on completed limits.
\end{proof}

\subsection{Compact rational quadrature}

\begin{definition}[Piecewise Lipschitz integrand]\label{def:lipschitz}
On a rational compact interval \([a,b]\), a piecewise Lipschitz presentation
contains rational breakpoints, a primitive-recursive rational evaluator
\(g(q,p)\) enclosing \(g(q)\) to precision \(2^{-p}\) at rational \(q\), and
rational constants \(L_j\) satisfying
\[
 \abs{g(x)-g(y)}\le L_j\abs{x-y}
\]
on each piece.  Endpoint values are supplied as one-sided names when needed.
\end{definition}

\begin{theorem}[Effective compact integration]\label{thm:compact-integral}
A piecewise Lipschitz presentation has a regular Riemann integral.  On one
piece of length \(T\), a mesh of width \(h\) gives a left-sum error at most
\(LTh\), plus the accumulated evaluation interval error.  Therefore a
primitive-recursive mesh size achieving error \(2^{-p}\) can be chosen.
\end{theorem}

\begin{proof}
For each cell \([x_k,x_{k+1}]\),
\[
 \abs{\int_{x_k}^{x_{k+1}}g(t)\,dt-g(x_k)(x_{k+1}-x_k)}
 \le\int_{x_k}^{x_{k+1}}L\abs{t-x_k}\,dt
 \le Lh(x_{k+1}-x_k).
\]
Summing yields \(LTh\).  A fully arithmetic proof replaces the integral in
this estimate by upper and lower rectangle sums: Lipschitz continuity bounds
the oscillation on a cell by \(Lh\), so the difference of Darboux sums is at
most \(LTh\).  The finite sums are rational and exact.  Split at the finitely
many breakpoints and allocate the error budget among pieces.
\end{proof}

\begin{definition}[Improper integral with an absolute-tail modulus]
\label{def:improper}
A presentation of \(g:\R\to\R\) includes compact piecewise-Lipschitz data and a
primitive-recursive radius \(R(p)\) satisfying
\[
 \int_{\abs x>R(p)}\abs{g(x)}\,dx\le2^{-p}.
\]
Its integral is the regular limit of compact integrals over
\([-R(p),R(p)]\).
\end{definition}

\begin{theorem}[Effective improper integration]\label{thm:improper}
The construction in \Cref{def:improper} produces a regular real independent of
the certified exhaustion.  Addition, scalar multiplication, and effective
changes of variables preserve it with computable moduli.
\end{theorem}

\begin{proof}
Two truncations beyond \(R(p)\) differ by an integral over a subset of the
absolute tail and hence by at most \(2^{-p}\).  Compact quadrature approximates
each truncation.  For addition, split the tail budget; for a change of variable,
transform the compact quadrature and the tail certificate.  Every output
precision therefore reduces to finitely many rational inequalities.
\end{proof}

\subsection{Densities and expectation}

\begin{definition}[Effective density]\label{def:density}
An effective density is a nonnegative integrand \(f\) with an effective
improper integral equal to one.  A function \(g\) is integrable against \(f\)
when \(\abs g f\) has an absolute-tail modulus and a compact quadrature
presentation.  Define
\[
 \E_f[g]=\int_{-\infty}^{\infty}g(x)f(x)\,dx.
\]
The random variable \(X(x)=x\) has mean \(\int xf(x)\,dx\) when
\(\abs x f(x)\) is effectively integrable.
\end{definition}

\begin{theorem}[Agreement with simple-law expectation]\label{thm:density-agreement}
Let \(f\) be an effective density and let \(g\) be effectively integrable
against \(f\).  There is a primitive-recursive sequence of finite
\emph{rational} laws whose values are rational cell representatives for
\(g\), whose successive levels admit certified finite couplings, and whose
\(L^1\) limit has expectation
\[
 \int_{-\infty}^{\infty}g(x)f(x)\,dx.
\]
Thus the density definition and the simple-law definition of expectation agree.
\end{theorem}

\begin{proof}
Fix an output precision \(p\).  Choose a rational radius \(R\) so that both
\(\int_{|x|>R}f\) and \(\int_{|x|>R}|g|f\) are below a small allocated
budget \(\eta\).  On \([-R,R]\), compute a rational bound \(G\) for
\(|g|\), and choose a rational partition so fine that the oscillation of
\(g\) on every cell is at most \(\eta\).  These choices are primitive
recursive from the compact-continuity and tail moduli.

For each cell \(I_j\), compact quadrature gives rational bounds
\(0\le l_j\le\int_{I_j}f\le u_j\) with
\(u_j-l_j\le\eta/J\), where \(J\) is the number of cells.  Put
\(q_j=l_j\) and add one residual atom of rational mass
\(q_*=1-\sum_jq_j\ge0\), represented at value zero.  The resulting masses
are nonnegative rationals summing exactly to one; a common denominator turns
them into natural weights.  Moreover,
\[
 q_*\le \int_{|x|>R}f+\sum_j(u_j-l_j),
\]
so the residual mass is explicitly small.  If \(r_j\in I_j\) is rational,
the expectation of the finite law is \(\sum_jq_jg(r_j)\).  Its difference
from \(\int gf\) is bounded by three finite estimates: cell oscillation,
\(G\sum_j(u_j-l_j)\), and the certified tail \(\int_{|x|>R}|g|f\).
Choosing \(\eta\) with guard bits makes their sum at most \(2^{-p}\).

For coherence between two levels, refine both rational partitions to their
finite common grid.  Match the rational masses greedily inside common cells;
the unmatched mass is bounded by the two quadrature deficits and the two tail
masses.  Transporting matched mass within a cell costs at most the cell
oscillation, while the unmatched compact mass costs at most \(2G\) times its
mass; the tail contribution is controlled by the \(|g|f\) modulus.  This
finite greedy transport is primitive recursive and supplies a rational
coupling with the required \(L^1\) bound.  Hence the levels form an effective
simple-law name.  The expectation estimate above and
\Cref{thm:l1-expect} show that its regular expectation is exactly the improper
integral.
\end{proof}

\begin{theorem}[Effective continuous LOTUS]\label{thm:continuous-lotus}
For an effective density \(f\) and an effectively integrable \(g\),
\[
 \E[g(X)]=\int_{-\infty}^{\infty}g(x)f(x)\,dx.
\]
For a primitive-recursive effective transformation \(Y=h(X)\), the expectation
of \(k(Y)\) is obtained by integrating \(k(h(x))f(x)\), without first
constructing a separate density for \(Y\).
\end{theorem}

\begin{proof}
The first statement is the definition justified by
\Cref{thm:density-agreement}.  For the second, apply the same finite simple-law
approximants with the composed evaluator \(k\circ h\).  This is the continuous
analogue of finite LOTUS and avoids an unnecessary inverse-function theorem.
\end{proof}

\subsection{CDF, Stieltjes, and layer-cake formulas}

\begin{definition}[Effective CDF name]\label{def:cdf}
For a density presentation, define
\[
 F(t)=\int_{-\infty}^{t}f(x)\,dx
\]
with a two-argument modulus controlling the left tail and compact quadrature.
More generally, a CDF/Stieltjes name consists of rational polygonal
approximants with monotonicity, endpoint, and variation certificates.
\end{definition}

\begin{theorem}[Effective integration by parts]\label{thm:stieltjes}
Let \(F\) be a certified CDF/Stieltjes name and let \(g\) be a compactly
supported piecewise \(C^1\) name, or a piecewise linear rational function.
Then
\[
 \int g\,dF=gF\big|_{a}^{b}-\int_a^bF(t)g'(t)\,dt
\]
with a primitive-recursive error modulus.  Passing to improper limits is valid
when the boundary terms and absolute tails have supplied moduli.
\end{theorem}

\begin{proof}
For polygonal approximants, the formula is finite summation by parts:
\[
 \sum_k g_k(F_{k+1}-F_k)
 =g_bF_b-g_aF_a-\sum_kF_{k+1}(g_{k+1}-g_k),
\]
up to a fixed endpoint convention.  This is a finite algebraic identity.
Uniform variation and continuity bounds control replacement of the names by
finer approximants.  The improper statement follows from the tail certificates.
\end{proof}

\begin{theorem}[Nonnegative layer cake]\label{thm:layer-cake}
For an effectively integrable nonnegative random variable,
\[
 \E[X]=\int_0^\infty\Pp(X>t)\,dt.
\]
For an integrable signed variable and any rational \(a\),
\[
 \E[X]-a
 =\int_a^\infty\Pp(X>t)\,dt
  -\int_{-\infty}^{a}\Pp(X\le t)\,dt.
\]
At \(a=\E[X]\), the two areas are equal.
\end{theorem}

\begin{proof}
At a finite rational simple level with denominator \(D\), the first formula is
exactly \Cref{thm:finite-layer-cake}: the integral is a sum of rectangles of
width \(1/D\).  Truncate the variable at a rational height and apply the
finite identity.  The integrability modulus bounds the discarded
\(\E[(X-B)^+]\), which is exactly the corresponding tail area at the simple
level.  Pass to the effective limit.

For signed \(X\), apply the first formula to \((X-a)^+\) and \((a-X)^+\), then
use \(X-a=(X-a)^+-(a-X)^+\).  The events defining the two positive parts are
\(X>t\) above \(a\) and \(X\le t\) below \(a\), with endpoint differences
confined to zero-width integration slices.  Setting \(a=\E[X]\) makes the left
side zero.
\end{proof}

\subsection{Absolute continuity: what is and is not internalized}

Classically, several formulations of absolute continuity are equivalent, and a
Radon--Nikodym theorem produces a density.  A bare existence conclusion is not
computational data.  We therefore use the following witnessed direction.

\begin{proposition}[Consequences of an explicit density witness]
\label{prop:ac-forward}
Given a nonnegative effective density \(f\), the set function
\(\mu([a,b])=\int_a^bf\) is effectively absolutely continuous with respect to
length on rational finite unions of intervals.  Its CDF is an effectively
absolutely continuous regular function, and its almost-everywhere derivative
is represented by the supplied \(f\) at all certified continuity points.
\end{proposition}

\begin{proof}
Given \(\varepsilon=2^{-p}\), first use the improper-tail modulus to choose
\(R\) with \(\int_{|x|>R}f<\varepsilon/2\).  The piecewise-Lipschitz data
on the compact interval \([-R,R]\) compute a rational upper bound \(M>0\)
for \(f\) there: evaluate one rational point on each piece and add its
Lipschitz constant times the piece length, with one guard bit for the evaluator.
If a finite union of intervals has total length below
\(\delta=\varepsilon/(2M)\), then its integral inside \([-R,R]\) is below
\(\varepsilon/2\), while its part outside contributes below the certified
spatial tail \(\varepsilon/2\).  This proves effective absolute continuity
without adding an independent amplitude-tail oracle.  The CDF increment
formula is its defining integral.  At a continuity point with a modulus, the
difference quotient is trapped between local lower and upper values of
\(f\), giving the derivative estimate.
\end{proof}

\begin{remark}[No uniform constructive Radon--Nikodym converse]
From a code that merely asserts an epsilon--delta property, one cannot in
general extract a primitive-recursive density evaluator and all its integration
moduli.  This report therefore does not claim the classical converse without a
witness.  If a density name is supplied as part of the certificate, all usual
forward equivalences become formalizable.
\end{remark}

\subsection{Cauchy truncation ambiguity}

\begin{proposition}[Cauchy expectation is undefined]\label{prop:cauchy-trunc}
For the standard Cauchy density
\(f(x)=1/(\pi(1+x^2))\), the positive and negative first moments both diverge.
The symmetric truncated first moment is zero, whereas truncation over
\([-R,2R]\) tends to \((\log2)/\pi\).  Thus a principal-value convention does
not define the expectation.
\end{proposition}

\begin{proof}
On a finite interval,
\[
 \int \frac{x}{\pi(1+x^2)}\,dx
 =\frac{1}{2\pi}\log(1+x^2).
\]
Symmetric endpoints cancel.  For \([-R,2R]\), the value is
\[
 \frac{1}{2\pi}\log\frac{1+4R^2}{1+R^2}
 \longrightarrow\frac{1}{2\pi}\log4=\frac{\log2}{\pi},
\]
with an explicit rational bound on the logarithmic ratio error.  On the
positive half-line,
\(x/(1+x^2)\ge1/(2x)\) for \(x\ge1\), so the positive moment dominates a
logarithmically divergent integral; symmetry gives divergence of the negative
part.  By \Cref{def:signed-extended}, the expectation is undefined.
\end{proof}

\section{Expected values of common distributions}\label{sec:distributions}

This section follows the article's table.  Finite distributions are exact
\(\ipra\) theorems.  Infinite and continuous distributions are effective
schemata whose parameters are rational or regular-real names with the displayed
positivity certificates.

\begin{proposition}[Bernoulli]\label{prop:bernoulli}
If \(X\in\{0,1\}\) with \(\Pp(X=1)=p\), then
\[
 \E[X]=p.
\]
\end{proposition}

\begin{proof}
The two weighted terms are \(0(1-p)\) and \(1p\).  In a natural-weight
presentation with weights \(W-r,r\), the equality is
\((0(W-r)+r)/W=r/W\).
\end{proof}

\begin{proposition}[Binomial]\label{prop:binomial}
For \(X\sim\operatorname{Bin}(n,p)\),
\[
 \E[X]=np.
\]
\end{proposition}

\begin{proof}
Construct \(X\) as the sum of \(n\) independent Bernoulli coordinates and use
finite linearity; independence is not actually needed for the mean once the
marginals are fixed.  Directly from the mass function,
\[
 \sum_{k=0}^{n}k\binom nkp^k(1-p)^{n-k}
 =np\sum_{k=1}^{n}\binom{n-1}{k-1}p^{k-1}(1-p)^{n-k}=np,
\]
using \(k\binom nk=n\binom{n-1}{k-1}\) and the finite binomial theorem.  For
rational \(p\), clear denominators; both combinatorial identities are proved by
primitive recursion.
\end{proof}

\begin{proposition}[Poisson]\label{prop:poisson}
For an effective nonnegative parameter \(\lambda\), the Poisson law
\[
 p_k=e^{-\lambda}\frac{\lambda^k}{k!}
\]
has expectation \(\lambda\).
\end{proposition}

\begin{proof}
At finite stage,
\[
 \sum_{k=0}^{N}kp_k
 =\lambda e^{-\lambda}\sum_{j=0}^{N-1}\frac{\lambda^j}{j!}.
\]
This is the exact shift identity \(k/k!=1/(k-1)!\).  Taylor's theorem with a
rational remainder bound gives a regular name for \(e^\lambda\) and proves
\(e^{-\lambda}\sum_{j\ge0}\lambda^j/j!=1\).  For a rational upper bound
\(L\ge\lambda\), choose \(N>2L\).  Successive absolute tail terms then have
ratio at most \(1/2\), so the tail is bounded by twice its first term.  This is
a primitive-recursive modulus.  Passing to the regular limit yields
\(\E[X]=\lambda\).
\end{proof}

\begin{proposition}[Geometric]\label{prop:geometric}
For rational \(0<p\le1\), let
\(\Pp(X=k)=p(1-p)^{k-1}\), \(k\ge1\).  Then
\[
 \E[X]=\frac1p.
\]
\end{proposition}

\begin{proof}
Put \(q=1-p\).  The finite polynomial identity
\[
 \sum_{k=1}^{N}kq^{k-1}
 =\frac{1-(N+1)q^N+Nq^{N+1}}{(1-q)^2}
\]
is proved by multiplying the sum by \(1-q\) twice and telescoping, with no
differentiation needed.  Multiplication by \(p=1-q\) gives a partial
expectation tending to \(1/p\).  Since \(q<1\), a rational separation
\(q\le1-1/d\) is computable from the rational code for \(p>0\).  The
finite binomial inequality \((1-1/d)^d\le1/2\) gives
\(q^{dm}\le2^{-m}\).  Bounded search therefore chooses \(m\) with
\((dm+d)2^{-m}\le2^{-r}\), yielding a primitive-recursive stage
\(N=dm\) at which \((N+1)q^N\le2^{-r}\).  This is an explicit tail
modulus; the case \(p=1\) has \(q=0\) and is immediate.
\end{proof}

\begin{proposition}[Uniform]\label{prop:uniform}
For rational \(a<b\), the uniform law on \([a,b]\) has
\[
 \E[X]=\frac{a+b}{2}.
\]
\end{proposition}

\begin{proof}
By exact polynomial integration,
\[
 \frac1{b-a}\int_a^b x\,dx
 =\frac{b^2-a^2}{2(b-a)}=\frac{a+b}{2}.
\]
Alternatively, the transformation \(x\mapsto a+b-x\) preserves the uniform
quadrature and pairs values with sum \(a+b\), forcing the mean to be its half.
The polynomial integral is already exact on every rational mesh.
\end{proof}

\begin{proposition}[Exponential]\label{prop:exponential}
For an effective positive parameter \(\lambda\), the density
\(f(x)=\lambda e^{-\lambda x}\) on \(x\ge0\) has
\[
 \E[X]=\frac1\lambda.
\]
\end{proposition}

\begin{proof}
On \([0,R]\), finite integration by parts gives
\[
 \int_0^R\lambda xe^{-\lambda x}\,dx
 =\frac1\lambda-e^{-\lambda R}\left(R+\frac1\lambda\right).
\]
The exponential Taylor bounds provide compact quadrature.  A positive rational
lower bound \(c\le\lambda\) gives
\(e^{-\lambda R}\le e^{-cR}\), and a primitive-recursive choice of \(R(p)\)
forces the displayed remainder below \(2^{-p}\).  Thus the improper integral
has regular value \(1/\lambda\).
\end{proof}

\begin{proposition}[Normal and standard normal]\label{prop:normal}
For an effective normal density with parameters \(\mu\) and \(\sigma>0\),
\[
 \E[X]=\mu.
\]
In particular, the standard normal mean is zero.
\end{proposition}

\begin{proof}
Substitute \(x=\mu+\sigma z\).  The expectation becomes
\[
 \mu\int_{-\infty}^{\infty}\varphi(z)\,dz
 +\sigma\int_{-\infty}^{\infty}z\varphi(z)\,dz,
\]
where \(\varphi(z)=(2\pi)^{-1/2}e^{-z^2/2}\).  The first integral is one by
the density certificate.  On every symmetric rational interval the second
integral is exactly zero because its integrand is odd and paired quadrature
nodes cancel.  A Gaussian tail certificate bounds the omitted absolute first
moment, so the improper limit is zero.  The transformation and all constants
are handled as regular interval names.
\end{proof}

\begin{proposition}[Pareto]\label{prop:pareto}
Let \(k>0\), \(\alpha>0\), and
\(f(x)=\alpha k^\alpha x^{-\alpha-1}\) for \(x\ge k\).  Then
\[
 \E[X]=
 \begin{cases}
 \dfrac{\alpha k}{\alpha-1},&\alpha>1,\\[4pt]
 +\infty,&0<\alpha\le1.
 \end{cases}
\]
Here rational powers are interpreted by regular names with monotonicity and
multiplication certificates.
\end{proposition}

\begin{proof}
The truncated mean is
\[
 \alpha k^\alpha\int_k^R x^{-\alpha}\,dx.
\]
For \(\alpha\ne1\), exact differentiation of the regular power name gives
\[
 \frac{\alpha k^\alpha}{1-\alpha}
   (R^{1-\alpha}-k^{1-\alpha}).
\]
If \(\alpha>1\), the first term tends effectively to zero and the remainder
simplifies to \(\alpha k/(\alpha-1)\).  If \(0<\alpha<1\), the positive term
\(R^{1-\alpha}\) has an explicit divergence modulus.  If \(\alpha=1\), the
truncated mean is \(k\log(R/k)\), which diverges by doubling \(R\).  All case
conditions are decidable for rational parameters; for regular parameters, a
strict separation certificate from one is required for the finite branch.
\end{proof}

\begin{proposition}[Cauchy]\label{prop:cauchy}
For a Cauchy law with location \(x_0\) and scale \(\gamma>0\), the expected
value is undefined.
\end{proposition}

\begin{proof}
Translation does not repair absolute integrability.  On the positive and
negative tails, \(\abs x\gamma/((x-x_0)^2+\gamma^2)\) is bounded below by a
positive multiple of \(1/\abs x\) beyond an explicitly computable radius.
Therefore both part expectations diverge logarithmically.  The signed
extended expectation is the indeterminate \(+\infty-\infty\), as in
\Cref{prop:cauchy-trunc}.
\end{proof}

\subsection{A mixed atom--density model}

\begin{proposition}[Precipitation occurrence and amount]
\label{prop:precipitation}
Suppose rain occurs with probability \(\alpha\), and conditional on rain the
amount is exponential with rate \(\lambda>0\); otherwise the amount is zero.
Then
\[
 \E[X]=\frac{\alpha}{\lambda}.
\]
\end{proposition}

\begin{proof}
Represent the model as a two-cell finite partition: dry with value zero and
wet with an effective exponential law.  The law of total expectation gives
\[
 \E[X]=(1-\alpha)\cdot0
 +\alpha\E[X\mid\text{wet}]
 =\alpha/\lambda.
\]
This is a direct combination of exact finite conditioning and the effective
exponential computation.
\end{proof}

\section{Convergence and interchange of expectation}\label{sec:convergence}

\subsection{Why pointwise convergence is insufficient}

\begin{example}[Moving spike]\label{ex:moving-spike}
Let the countable sample space be \(\N\) with
\(\Pp(\{k\})=2^{-k-1}\).  Define
\[
 X_n(k)=2^{n+1}\ind_{k=n}.
\]
For every fixed \(k\), \(X_n(k)=0\) for all \(n>k\), so \(X_n\to0\)
pointwise.  Nevertheless,
\[
 \E[X_n]=2^{-n-1}2^{n+1}=1
\]
for every \(n\).  Hence pointwise convergence alone does not imply convergence
of expectations.
\end{example}

This counterexample is fully effective: the geometric law has an exact tail,
the values are natural numbers, and the pointwise convergence modulus at
outcome \(k\) is \(k+1\).  What fails is a uniform integrability bound.

\begin{example}[Interchanging an infinite sum and expectation can fail]
\label{ex:series-interchange}
Put \(X_{-1}=0\), \(Y_0=X_0\), and
\(Y_n=X_n-X_{n-1}\) for \(n\ge1\), using the moving spikes above.  Then
\[
 \sum_{j=0}^{N}Y_j=X_N\longrightarrow0
\]
pointwise, so the pointwise series sum is zero.  But
\[
 \sum_{j=0}^{N}\E[Y_j]=\E[X_N]=1
\]
for every \(N\).  Thus
\(\E[\sum_jY_j]=0\ne1=\sum_j\E[Y_j]\).
Absolute summability or domination hypotheses are essential.
\end{example}

\subsection{Effective monotone convergence}

\begin{theorem}[Monotone convergence as a lower-real identity]
\label{thm:monotone-convergence}
Let \(0\le X_0\le X_1\le\cdots\) be a primitive-recursive sequence of
nonnegative finite simple-law names on compatible refinements.  Define the
pointwise supremum by its increasing approximants.  Then
\[
 \E[\sup_nX_n]=\sup_n\E[X_n]
\]
as equality of lower reals.  If an expectation-tail modulus is supplied, both
sides are equal as regular reals.
\end{theorem}

\begin{proof}
At every finite stage, expectation preserves order, so
\(\E[X_n]\) is increasing and bounded above by every finite approximation to
\(\E[\sup X_n]\).  Conversely, the definition of the supremum name uses
\(X_n\) itself as its increasing simple approximants; therefore its lower-real
expectation is, by construction, the directed family of the finite
expectations.  More explicitly, any rational lower bound certified for one
side occurs at a finite level and is certified for the other side at the same
level.  If a modulus says the omitted expectation after stage \(N(p)\) is at
most \(2^{-p}\), the common lower real becomes a regular interval name.
\end{proof}

\begin{corollary}[Nonnegative series and Tonelli form]
\label{cor:tonelli}
For a primitive-recursive sequence of nonnegative observables,
\[
 \E\left[\sum_{n=0}^{\infty}X_n\right]
 =\sum_{n=0}^{\infty}\E[X_n]
\]
as lower reals.  With an absolute expectation-tail modulus, the equality is
one of regular reals and the sum may be effectively regrouped.
\end{corollary}

\begin{proof}
Apply \Cref{thm:monotone-convergence} to the finite partial sums and use exact
finite linearity at every stage.
\end{proof}

\subsection{Finite Fatou kernel and effective liminf}

\begin{lemma}[Finite Fatou kernel]\label{lem:finite-fatou}
For nonnegative finite observables \(X_N,\ldots,X_M\),
\[
 \E\left[\min_{N\le n\le M}X_n\right]
 \le \min_{N\le n\le M}\E[X_n].
\]
\end{lemma}

\begin{proof}
The pointwise minimum is at most each \(X_n\).  Monotonicity gives an
expectation inequality for each \(n\); bounded search takes the least rational
right side.  The minimum evaluator and the selected index are primitive
recursive on the finite list.
\end{proof}

\begin{theorem}[Effective Fatou lemma]\label{thm:fatou}
Let \(X_n\ge0\) be an effective sequence.  For every \(N\), suppose the
finite minima
\[
 Y_{N,M}=\min_{N\le n\le M}X_n
\]
are supplied on compatible simple-law refinements and converge, as
\(M\to\infty\), to a name \(Y_N\) with an explicit \(L^1\) modulus.  Suppose
also that the increasing sequence \(Y_N\) presents the pointwise
\(\liminf\).  Then
\[
 \E[\liminfN_{n\to\infty}X_n]
 \le\liminfN_{n\to\infty}\E[X_n]
\]
as a lower-real inequality; with regular moduli it is an inequality of regular
reals.
\end{theorem}

\begin{proof}
For finite \(M\), \Cref{lem:finite-fatou} gives
\[
 \E[Y_{N,M}]
 \le\min_{N\le n\le M}\E[X_n].
\]
The supplied \(L^1\) modulus and
\(\lvert\E U-\E V\rvert\le\E\lvert U-V\rvert\) let \(M\) increase on the
left, yielding \(\E[Y_N]\).  The finite minima on the right decrease to the
upper-real tail infimum \(\inf_{n\ge N}\E[X_n]\); a regular version uses the
corresponding numerical modulus.  Now \(Y_N\) increases with \(N\), so
effective monotone convergence gives the expectation of the liminf.  Taking
the increasing supremum over \(N\) on the right gives the numerical liminf.
At each requested precision only finite minima, one finite expectation theorem,
and the stated moduli are used.
\end{proof}

\subsection{Uniform integrability and dominated convergence}

\begin{definition}[Effective uniform integrability]\label{def:ui}
A sequence \((X_n)\) is effectively uniformly integrable when a
primitive-recursive threshold \(K(p)\) satisfies
\[
 \E\bigl[\abs{X_n}\ind_{\abs{X_n}>K(p)}\bigr]\le2^{-p}
 \quad\text{for every }n.
\]
The universal quantifier is implemented by a proof schema for the fixed
sequence evaluator, not by bounded search over all \(n\).
\end{definition}

\begin{lemma}[Truncation estimate]\label{lem:truncation}
Suppose \(\abs X,\abs Y\le G\), where \(G\ge0\).  For rational
\(K,\eta>0\),
\[
 \E\abs{X-Y}
 \le \eta+2K\Pp(\abs{X-Y}>\eta)
      +2\E[G\ind_{G>K}].
\]
\end{lemma}

\begin{proof}
Split the sample space into three decidable regions.  On
\(\abs{X-Y}\le\eta\), the contribution is at most \(\eta\).  On the remaining
region with \(G\le K\), we have \(\abs{X-Y}\le\abs X+\abs Y\le2K\), giving the
middle term.  On \(G>K\), use \(\abs{X-Y}\le2G\).  Indicator linearity and the
finite union bound produce the displayed inequality.  For effective names,
apply it on a sufficiently fine common finite refinement and absorb the
\(L^1\) approximation errors.
\end{proof}

\begin{theorem}[Effective dominated convergence]\label{thm:dct}
Let \(X_n,X,G\) be effective \(L^1\) names satisfying
\(\abs{X_n},\abs X\le G\).  Assume:
\begin{enumerate}
\item an integrability-tail modulus for \(G\), and
\item a convergence-in-probability modulus: given \(p,q\), a
  primitive-recursive \(N(p,q)\) makes
  \[
   n\ge N(p,q)\Longrightarrow
   \Pp(\abs{X_n-X}>2^{-p})\le2^{-q}.
  \]
\end{enumerate}
Then \(X_n\to X\) in \(L^1\) with a computable modulus, and consequently
\[
 \E[X_n]\longrightarrow\E[X].
\]
\end{theorem}

\begin{proof}
Given output precision \(p\), choose \(K\) so that
\(2\E[G\ind_{G>K}]\le2^{-p-2}\).  Put
\(\eta=2^{-p-2}\).  Choose \(q\) so that
\(2K2^{-q}\le2^{-p-2}\), and then choose
\(n\ge N(p+2,q)\).  The truncation estimate gives
\(\E\abs{X_n-X}\le3\cdot2^{-p-2}<2^{-p}\), after allowing one additional guard
bit for name approximation.  Finally,
\[
 \abs{\E[X_n]-\E[X]}\le\E\abs{X_n-X}\le2^{-p}.
\]
Every choice is primitive recursive from the supplied moduli and rational
bounds.
\end{proof}

\begin{corollary}[Bounded convergence]\label{cor:bounded-convergence}
If \(\abs{X_n},\abs X\le M\) for a fixed rational \(M\) and \(X_n\to X\) in
probability with a modulus, then expectations converge effectively.
\end{corollary}

\begin{proof}
Take \(G=M\).  Its tail above any \(K>M\) is zero, so
\Cref{thm:dct} applies.  Directly,
\(\E\abs{X_n-X}\le\eta+2M\Pp(\abs{X_n-X}>\eta)\).
\end{proof}

\begin{theorem}[Vitali form]\label{thm:vitali}
Effective convergence in probability plus effective uniform integrability of
\((X_n)\) and the limit \(X\) implies effective \(L^1\) convergence and hence
convergence of expectations.
\end{theorem}

\begin{proof}
Use \Cref{lem:truncation} with separate tail bounds for \(X_n\) and \(X\):
\[
 \abs{X_n-X}\le
 \abs{X_n-X}\ind_{\abs{X_n}\le K,\abs X\le K}
 +\abs{X_n}\ind_{\abs{X_n}>K}
 +\abs X\ind_{\abs X>K}.
\]
The first term is bounded by
\(\eta+2K\ind_{\abs{X_n-X}>\eta}\).  Choose \(K,\eta\), and the probability
precision exactly as in \Cref{thm:dct}.
\end{proof}

\subsection{Summary of constructive hypotheses}

The classical convergence theorems remain valid, but their useful
\(\ipra\) forms expose the missing finite data:
\begin{itemize}
\item monotone convergence needs no rate for a lower-real conclusion, but a
  regular-real value needs a tail modulus;
\item Fatou needs effective names for the finite minima and the limiting
  infima;
\item dominated convergence needs an integrable-tail modulus and a convergence
  modulus;
\item exchanging expectation with a signed infinite sum needs absolute
  expectation tails or another uniform-integrability certificate.
\end{itemize}
The moving-spike examples show that these are mathematical requirements, not
merely artifacts of formalization.

\section{Characteristic functions, Fourier formulas, and applications}\label{sec:fourier}

\subsection{Characteristic functions and the first moment}

For a finite rational law, define the regular complex name
\[
 \phi_X(t)=\E[e^{itX}]
 =\E[\cos(tX)]+i\E[\sin(tX)].
\]
The sine and cosine names are obtained from finite Taylor sums with explicit
remainders after range reduction.

\begin{theorem}[Derivative of the characteristic function at zero]
\label{thm:char-derivative}
If \(\E[X^2]\) is finite with an effective name, then for every nonzero
rational \(h\),
\[
 \left|\frac{\phi_X(h)-1}{h}-i\E[X]\right|
 \le\frac{\abs h}{2}\E[X^2].
\]
Consequently \(\phi_X'(0)=i\E[X]\) effectively.
\end{theorem}

\begin{proof}
The elementary Taylor remainder estimate
\[
 \abs{e^{iu}-1-iu}\le\frac{u^2}{2}
\]
is proved constructively from the integral or alternating-series remainder for
sine and cosine; at rational precision it is a finite polynomial inequality
plus a geometric tail bound.  Put \(u=hX_i\), divide by \(\abs h\), and take
expectations:
\begin{align*}
\left|\frac{\phi_X(h)-1}{h}-i\E[X]\right|
&=\left|\E\left[\frac{e^{ihX}-1-ihX}{h}\right]\right|\\
&\le\E\left[\frac{\abs{e^{ihX}-1-ihX}}{\abs h}\right]\\
&\le\frac{\abs h}{2}\E[X^2].
\end{align*}
For an effective infinite law, first truncate using the second-moment tail
modulus.  Given output precision, choose \(h\) small enough that the final
bound is below it.
\end{proof}

\subsection{Finite cyclic Fourier inversion}

\begin{theorem}[Finite Fourier expectation identity]\label{thm:finite-fourier}
Let \(p_x\) be a rational probability mass on \(\Z/N\Z\), let
\(\omega=e^{2\pi i/N}\), and define
\[
 \phi(k)=\sum_{x<N}p_x\omega^{kx},\qquad
 \widehat g(k)=\sum_{x<N}g(x)\omega^{-kx}.
\]
Then
\[
 \E[g(X)]=\frac1N\sum_{k<N}\widehat g(k)\phi(k).
\]
This identity has an effective \(\ipra\) realization using regular complex
names, and an exact algebraic realization in a cyclotomic presentation.
\end{theorem}

\begin{proof}
Expand the right side and use finite Fubini:
\[
 \frac1N\sum_{x,y<N}g(x)p_y
       \sum_{k<N}\omega^{k(y-x)}.
\]
The inner sum is \(N\) if \(x=y\) and zero otherwise.  For \(x\ne y\), this
is the finite geometric identity
\[
 (1-\omega^{y-x})\sum_{k<N}\omega^{k(y-x)}
 =1-\omega^{N(y-x)}=0,
\]
together with a primitive-root certificate that
\(1-\omega^{y-x}\ne0\).  Thus only diagonal terms remain, giving
\(\sum_xg(x)p_x\).  With regular complex names, approximate all roots of unity
with guard bits and use the same exact exponent arithmetic; with a cyclotomic
algebra code, the orthogonality certificate is symbolic.
\end{proof}

\subsection{Effective continuous Fourier formula}

\begin{theorem}[Expectation through Fourier inversion]
\label{thm:continuous-fourier}
Suppose a density \(f\), its characteristic function \(\phi\), a function
\(g\), and its Fourier transform
\[
 G(t)=\int_{\R}g(x)e^{-itx}\,dx
\]
are given by effective integration names.  Assume explicit absolute-tail and
Fubini moduli for the double integral.  Then, with the article's sign
convention,
\[
 \E[g(X)]=\frac1{2\pi}\int_{\R}G(t)\phi(t)\,dt.
\]
\end{theorem}

\begin{proof}
Truncate both \(x\) and \(t\) to compact rational intervals chosen from the
absolute-tail modulus.  Approximate both integrals by rational rectangles.
On the resulting finite double sum, interchange the order by exact finite
Fubini and substitute the finite Fourier inversion approximation to \(f\).
The accumulated quadrature, inversion, and tail errors are allocated less than
\(2^{-p}/4\) each.  Since the moduli compute all truncation radii and mesh
sizes from \(p\), this is an effective theorem schema.  Without absolute
Fubini data, the order change is not claimed.
\end{proof}

\subsection{Moments and finite moment determinacy}

\begin{definition}[Moments]
The \(r\)-th raw moment is \(m_r=\E[X^r]\), and the \(r\)-th central moment is
\(\E[(X-\E X)^r]\).  For finite rational laws these are exact rational
expectations.  Their finite lists are primitive recursive in the law code.
\end{definition}

\begin{theorem}[Finite-support moment determinacy]\label{thm:moment-determinacy}
Let \(X\) be supported on \(\{0,1,\ldots,B\}\).  The moments
\(m_0,m_1,\ldots,m_B\) determine the probability masses uniquely, and the
masses are primitive-recursively recoverable by rational arithmetic.
\end{theorem}

\begin{proof}
Use falling factorials
\((x)_r=x(x-1)\cdots(x-r+1)\) and factorial moments
\(f_r=\E[(X)_r]\).  Ordinary powers and falling factorials are related by the
finite triangular Stirling-number transformations, whose integer coefficients
are primitive recursive.  Hence \((m_0,\ldots,m_B)\) determines
\((f_0,\ldots,f_B)\).

If \(p_k=\Pp(X=k)\), then
\[
 f_r=\sum_{k=r}^{B}(k)_r p_k.
\]
This is upper triangular, with diagonal \((r)_r=r!>0\).  Recover the masses by
back substitution:
\[
 p_B=\frac{f_B}{B!},\qquad
 p_r=\frac{f_r-\sum_{k=r+1}^{B}(k)_r p_k}{r!}.
\]
A downward bounded recursion computes every \(p_r\).  Uniqueness follows from
the same recursion applied to two candidate vectors.  This is a complete
\(\ipra\) proof of moment determinacy for known finite support; no analytic
moment problem is involved.
\end{proof}

\begin{corollary}[Finite moment generating function]
For finite support, the formal polynomial/exponential sum
\(M(t)=\E[e^{tX}]\) determines the law, since its first \(B\) derivatives at
zero give the moments in \Cref{thm:moment-determinacy}.
\end{corollary}

\subsection{Expected utility and finite decision problems}

\begin{theorem}[Finite expected-utility maximization]\label{thm:expected-utility}
Let \(A\) be a finite set of actions, \(S\) a finite weighted state law, and
\(u(a,s)\in\Q\) a primitive-recursive utility table.  There is a
primitive-recursively selected action \(a^*\) satisfying
\[
 \E[u(a^*,S)]\ge\E[u(a,S)]\quad\text{for every }a\in A.
\]
\end{theorem}

\begin{proof}
Compute the exact rational expected utility for every action.  Bounded search
selects the least index attaining the maximum, using decidable rational order.
An induction over the action list proves that the running candidate dominates
all earlier candidates.  The output includes its comparison certificate.
\end{proof}

\begin{remark}
This proves the computational content of finite expected-utility choice.  It
does not assert that every preference ordering has a utility representation;
that is a separate representation theorem with its own hypotheses.
\end{remark}

\subsection{Indicators, Monte Carlo, and empirical means}

\begin{theorem}[Probability as expectation and unbiased frequency]
\label{thm:monte-carlo}
For an event \(A\), \(\Pp(A)=\E[\ind_A]\).  For iid repetitions and empirical
frequency
\[
 \widehat p_n=\frac1n\sum_{j<n}\ind_A^{(j)},
\]
we have
\[
 \E[\widehat p_n]=\Pp(A),\qquad
 \Var(\widehat p_n)=\frac{p(1-p)}n.
\]
Consequently Chebyshev supplies the exact finite Monte Carlo error contract
\[
 \Pp(\abs{\widehat p_n-p}\ge\varepsilon)
 \le\frac{p(1-p)}{n\varepsilon^2}
 \le\frac1{4n\varepsilon^2}.
\]
\end{theorem}

\begin{proof}
The indicator identity was proved in \Cref{thm:indicator}.  Apply the
sample-mean and variance theorems to Bernoulli variables, whose variance is
\(p-p^2=p(1-p)\).  Finally,
\(p(1-p)\le1/4\) is the square identity
\(1/4-p(1-p)=(p-1/2)^2\ge0\), interpreted after clearing denominators.
\end{proof}

\subsection{Center of mass}

\begin{theorem}[Balance point]\label{thm:center-mass}
Place nonnegative rational masses \(m_i\) at rational positions \(x_i\), with
positive total mass \(M=\sum_i m_i\).  The unique point \(c\) with zero total
torque
\[
 \sum_i m_i(x_i-c)=0
\]
is
\[
 c=\frac{\sum_i m_ix_i}{M}.
\]
Thus the center of mass is exactly the expectation under probabilities
\(p_i=m_i/M\).
\end{theorem}

\begin{proof}
Expand the torque:
\[
 \sum_i m_ix_i-c\sum_i m_i=\sum_i m_ix_i-cM.
\]
It vanishes at the displayed \(c\).  Conversely, if it vanishes at \(d\), then
\(dM=\sum_i m_ix_i=cM\); positive \(M\) permits cancellation.  This is a
single denominator-cleared linear identity.
\end{proof}

\begin{proposition}[Beta balance example]\label{prop:beta}
For positive natural \(\alpha,\beta\), the beta density proportional to
\(x^{\alpha-1}(1-x)^{\beta-1}\) on \([0,1]\) has mean
\(\alpha/(\alpha+\beta)\).
\end{proposition}

\begin{proof}
For natural parameters, the density is a polynomial times a rational
normalization.  Integration by parts gives the finite beta recursion
\[
 B(\alpha+1,\beta)=\frac{\alpha}{\alpha+\beta}B(\alpha,\beta).
\]
The expectation is the ratio of these two exact polynomial integrals.  The
recursion can alternatively be proved by expanding
\((1-x)^{\beta-1}\) and integrating monomials rationally.
\end{proof}

\subsection{A bounded Wald identity}

\begin{theorem}[Wald equation for a bounded stopping time]
\label{thm:wald}
Let \(X_0,\ldots,X_{N-1}\) be iid finite rational variables with mean \(\mu\).
Let \(T\le N\) be a stopping time, meaning the event \(T>j\) depends only on
\(X_0,\ldots,X_{j-1}\).  Then
\[
 \E\left[\sum_{j<T}X_j\right]=\mu\E[T].
\]
\end{theorem}

\begin{proof}
Pointwise,
\[
 \sum_{j<T}X_j=\sum_{j<N}X_j\ind_{T>j},
 \qquad
 T=\sum_{j<N}\ind_{T>j}.
\]
Finite linearity reduces the first expectation to a sum.  By the stopping-time
condition, \(\ind_{T>j}\) is a function of earlier coordinates and is
independent of \(X_j\).  Therefore
\[
 \E[X_j\ind_{T>j}]=\E[X_j]\Pp(T>j)=\mu\Pp(T>j).
\]
Sum over \(j<N\) and use the tail formula for the bounded integer \(T\).
\end{proof}

\subsection{Finite-dimensional quantum expectation}

\begin{theorem}[Quantum expectation and variance identity]
\label{thm:quantum}
Let \(\psi\in\Q(i)^d\) satisfy \(\psi^*\psi=1\), and let
\(A\in\Q(i)^{d\times d}\) be Hermitian.  Put
\[
 \langle A\rangle=\psi^*A\psi.
\]
Then \(\langle A\rangle\) is real and
\[
 (\Delta A)^2
 \defeq\psi^*A^2\psi-(\psi^*A\psi)^2
 =\norm{(A-\langle A\rangle I)\psi}^2\ge0.
\]
\end{theorem}

\begin{proof}
Hermiticity gives
\((\psi^*A\psi)^*=\psi^*A^*\psi=\psi^*A\psi\), so the scalar is real.  Expand
the squared norm:
\begin{align*}
\psi^*(A-\mu I)^*(A-\mu I)\psi
 &=\psi^*(A^2-2\mu A+\mu^2I)\psi\\
 &=\psi^*A^2\psi-2\mu^2+\mu^2
 =\psi^*A^2\psi-\mu^2,
\end{align*}
where \(\mu=\psi^*A\psi\) and \(\psi^*\psi=1\).  The norm square is a finite
sum of rational squares of real and imaginary parts, hence nonnegative in
\(\ipra\).  Algebraic or regular complex entries are handled by the same
finite identity at each precision.
\end{proof}

\section{Lean formalization and object-language route}\label{sec:lean}

\subsection{Repository interfaces reused}

The ProveIt repository has two distinct levels that must not be conflated.
The ambient Lean metatheory provides convenient finite sums, rationals, and
proof automation.  The internal target is an explicit derivation
\[
 ([]:\operatorname{List}(\mathrm{Fm}\;\mathrm{PTm}\;k))
 \;\vdash[\mathrm{PRAAx}]\;\varphi
\]
in the formalized \(\ipra\) calculus.  The existing
\code{pra_iff_native} theorem then transports a closed \code{PRAAx} derivation
to the concise native scalar core.

The checked arithmetic library already contains the principal combinators
needed here:
\begin{itemize}
\item \code{IPRADerivations.Core.BoundedSums}: elimination of terms from a
  zero nonnegative sum and construction of a zero sum from pointwise zeros;
\item \code{IPRADerivations.Core.SumLemmas}: sum shift, sum of pointwise sums,
  bounded congruence, and ``a term is at most a nonnegative sum'';
\item \code{IPRADerivations.Core.MultiSums}: rectangular finite Fubini and
  reindexing;
\item \code{IPRADerivations.Core.Pairing}: coding product outcomes;
\item arithmetic, order, truncated subtraction, division with remainder, and
  ring-normalization modules used by \code{deriveNat};
\item \code{ConciseIPRA.NativeBiInterpretation}: the
  \code{pra_iff_native} bridge.
\end{itemize}

The earlier constructive central-limit packet supplies compatible design
patterns: a natural-weight \code{FinLaw}, nested rational interval
\code{Regular} names, and an \code{InternalProof} structure separating a
primitive-recursive checker from its actual \code{PRAAx} derivation.  The
present packet follows those patterns but does not assume the uncompleted CLT
blueprints as theorems.

\subsection{Ambient finite API}

The file \code{lean/ExpectedValueFinite.lean} begins with

\begin{lstlisting}[language={}]
structure FinLaw where
  card : Nat
  weight : Fin card -> Nat
  total : Nat
  total_pos : 0 < total
  weight_sum : (sum i, weight i) = total

def FinLaw.expect (L : FinLaw)
    (X : Fin L.card -> Rat) : Rat :=
  (sum i, (L.weight i : Rat) * X i) / L.total
\end{lstlisting}

It also defines event probability, moments, variance, finite CDFs, almost-sure
equality, pushforward fiber weights, product-law data, partition cell weights,
and conditional means.  The structures
\code{FiniteExpectationProofPackage},
\code{FiniteDependenceProofPackage}, and \code{SampleProofPackage} list the
ambient theorems to be constructed.  They are records of obligations, not
axioms: no inhabitant is provided.

A good ambient implementation should avoid a single enormous proof.  The
recommended module graph is

\begin{verbatim}
IPRADerivations/Probability/FiniteLaw.lean
IPRADerivations/Probability/Expectation.lean
IPRADerivations/Probability/TailSum.lean
IPRADerivations/Probability/Variance.lean
IPRADerivations/Probability/MarkovChebyshev.lean
IPRADerivations/Probability/CauchySchwarz.lean
IPRADerivations/Probability/Independence.lean
IPRADerivations/Probability/Conditional.lean
IPRADerivations/Probability/SampleMean.lean
IPRADerivations/Probability/Kolmogorov.lean
\end{verbatim}

Each file should first prove the theorem over ordinary \code{Rat} and finite
sums.  More importantly, it should state a denominator-cleared lemma whose
shape already matches the natural-number object language.

\subsection{Natural encodings}

Introduce codes such as

\begin{lstlisting}[language={}]
structure ZCode where
  pos neg : Nat

structure QCode where
  num : ZCode
  den : Nat
  denPos : 0 < den

structure FinLawCode where
  card total commonDen : Nat
  weightCode valueCode : Nat
\end{lstlisting}

A finite list is encoded by the repository's pairing/list scheme.  Access
outside the declared length returns zero.  The decoder validates:
\begin{enumerate}
\item positive total and positive common denominator;
\item bounded weight sum equals the total;
\item every observable entry has the declared common denominator;
\item theorem-specific side conditions, such as nonnegativity or a positive
  threshold.
\end{enumerate}
Invalid inputs return zero from the checker.  Thus a universally quantified
open equation \code{check(args)=0} acts as an implication: malformed data pass
vacuously, while valid data trigger the mathematical verification.

The essential primitive-recursive functions are
\begin{verbatim}
weightSum       weightedNum      eventMass
absWeightedNum secondMomentNum  fiberMass
integerTailMass productIndex    cellMass
firstCrossing  finiteMoment     regularIntervalCheck
\end{verbatim}
Their defining equations must be proved once and reused.  No semantic theorem
should unfold a large decoder repeatedly.

\subsection{Object-language theorem shapes}

The file \code{lean/ObjectLanguageTargets.lean} defines a generic open target

\begin{lstlisting}[language={}]
def checkedOpen {k : Nat} (check : PR k) : Fm PTm k :=
  PTm.app check Syn.var = PTm.zero

structure InternalProof {k : Nat}
    (semantic : (Fin k -> Nat) -> Nat)
    extends Representation semantic where
  derived : ([] : List (Fm PTm k))
    |- [PRAAx] checkedOpen check
\end{lstlisting}

The concrete syntax in the source uses ProveIt's equality and derivability
symbols; the ASCII display above is only schematic.  A completed theorem
contains both:
\begin{enumerate}
\item \code{correct}: the \code{PR} term evaluates to the intended semantic
  test; and
\item \code{derived}: an actual derivation proving that the test always
  returns zero.
\end{enumerate}
Standard-model truth establishes neither field automatically.

\subsection{Proof plan for each finite theorem}

\paragraph{Linearity.}
Normalize \(X,Y\) to a common denominator.  The numerator for \(X+Y\) is the
bounded sum of the pointwise numerator sum.  Instantiate \code{sum_add}, apply
the ring tool to denominator products, and close universally.

\paragraph{Monotonicity.}
The checker stores nonnegative numerator differences
\(d_i=Y_i-X_i\).  A bounded sum of nonnegative naturals is nonnegative.
Equivalently, define the target as the existence of a natural slack
\(s\) with \(\operatorname{num}(Y)-\operatorname{num}(X)=s\); the sum of the
pointwise slacks supplies \(s\).

\paragraph{Indicator and finite LOTUS.}
Indicator expectation unfolds to \code{eventMass}.  For LOTUS, define
\code{fiberMass(v)} as a bounded sum.  Apply \code{MultiSums} to the matrix
\(w_i\ind_{f(i)=v}g(v)\), and use the proof that exactly one finite fiber
indicator is active for each \(i\).

\paragraph{Tail sum.}
First derive the open identity
\[
 x=\sum_{k<B}\ind_{k<x}\qquad(x\le B)
\]
by quantifier-free induction on \(x\) or \(B\).  Substitute each encoded value,
multiply by the weight, and use finite Fubini.

\paragraph{Variance.}
After common-denominator normalization, invoke ring normalization on the
identity
\[
 W\sum_iw_i(X_i-\mu)^2
 =W\sum_iw_iX_i^2-(\sum_iw_iX_i)^2.
\]
The displayed form avoids division and is polynomial over signed-natural
codes.

\paragraph{Markov and Chebyshev.}
Compile the pointwise slack in
\(a\ind_{a\le X_i}\le X_i\), sum it, and clear denominators.  Chebyshev is the
same theorem applied to the encoded square deviation.  There is no need to
formalize a general probability inequality framework before these two files.

\paragraph{Cauchy--Schwarz.}
Use the weighted Lagrange identity from \Cref{lem:lagrange}.  The double sum of
squares is encoded by \code{MultiSums}; nonnegativity is a direct natural
statement.  This is more proof-friendly than a minimization or quadratic-
discriminant argument.

\paragraph{Independence and conditioning.}
For product laws, mixed-radix pairing and \code{MultiSums} prove factorization.
For conditioning, split on the natural cell mass.  A positive cell uses exact
cancellation; a zero cell uses \code{sum_zero_elim} to show all its weighted
terms vanish.  Sum the cell identities for total expectation.

\paragraph{Kolmogorov.}
Encode the first-crossing index by bounded search through the partial sums.
Prove disjointness of first-crossing indicators and factor the tail increment
using the product-law theorem.  The remaining argument is finite sum algebra.
This is the largest finite proof, so it should be attempted only after all
previous modules compile.

\subsection{Regular expectation API}

The file \code{lean/RegularExpectation.lean} defines nested rational intervals,
finite simple approximation levels, refinement maps preserving mass, an exact
refinement \(L^1\) distance, absolute-series names, and compact/improper
integrand interfaces.  The next implementation modules should be

\begin{verbatim}
IPRADerivations/Probability/RegularExpectation.lean
IPRADerivations/Probability/Series.lean
IPRADerivations/Probability/Integration.lean
IPRADerivations/Probability/Distributions.lean
IPRADerivations/Probability/Convergence.lean
IPRADerivations/Probability/Fourier.lean
\end{verbatim}

At this level the object-language sentence always contains a precision
variable.  A checker decodes a presentation and verifies one interval
inclusion with dyadic slack.  The mathematical proof computes the finite level
needed for that precision, invokes the finite kernel, and combines rational
error bounds.  It does not quantify over an ambient \code{Real}.

\subsection{Suggested milestone sequence}

\begin{enumerate}
\item Compile the ambient finite file against the repository's pinned Lean and
  Mathlib version, replacing every proof-package field by an actual theorem.
\item Add natural encoders and prove decoder/encoder round trips for valid
  codes.
\item Construct object-language derivations for normalization, constants,
  additivity, homogeneity, and monotonicity.
\item Add indicators, LOTUS, tail sums, and variance.
\item Add Markov, Chebyshev, and Cauchy--Schwarz.
\item Add product laws, conditioning, sample means, and Kolmogorov.
\item Only then implement regular names, series, quadrature, convergence, and
  Fourier modules.
\item Universally close each theorem and transport it with
  \code{pra_iff_native}; aggregate packages only after individual modules are
  checked.
\end{enumerate}

\subsection{Completion criterion}

The formalization is complete only when all of the following hold:
\begin{itemize}
\item the ambient files compile with the repository's pinned toolchain;
\item every decoder and primitive-recursive evaluator has proved defining
  equations;
\item each main theorem contains an explicit \code{Prf PRAAx} derivation;
\item no theorem is introduced through \code{axiom}, \code{sorry}, or a
  constant-zero checker masquerading as the semantic verifier;
\item the native proof is obtained through the existing proved bridge;
\item the Python validation remains outside the trusted proof boundary.
\end{itemize}

\section{Coverage ledger, conclusions, and further research}\label{sec:conclusion}

\subsection{Coverage of the source article}

The accompanying \code{THEOREM\_LEDGER.md} contains 94 numbered mathematical
items extracted from the consulted revision of the article.  The coverage is
organized as follows.

\begin{longtable}{@{}p{0.22\textwidth}p{0.68\textwidth}@{}}
\toprule
Article region & Constructive treatment \\
\midrule
\endhead
Finite definition and examples & Exact natural-weight/rational-law
formalization; die and roulette computations; equiprobable arithmetic mean;
vector and matrix expectation. \\
Countable definition & Absolute-tail regular sums; effective reindexing;
conditional rearrangement counterexample; logarithmic law; extended
expectations; St Petersburg divergence; undefined two-sided example. \\
Density and general integral & Effective simple-law names, compact quadrature,
improper tails, density agreement, LOTUS, Stieltjes integration, and layer
cake.  Existence-only density extraction is replaced by a witness. \\
Common distributions & Bernoulli and binomial exactly; Poisson, geometric,
uniform, exponential, normal, Pareto, and Cauchy with explicit series or
integration moduli. \\
Basic properties & Nonnegativity, finite linearity, monotonicity,
nondegeneracy, almost-sure congruence, constants, absolute-value bound,
indicators, CDF/tail formulas, nonmultiplicativity, independence, and
multidimensional LOTUS. \\
Inequalities & Markov, Chebyshev, certified Jensen, Lyapunov, Holder,
Cauchy--Schwarz, and Minkowski. \\
Convergence & Effective monotone convergence, nonnegative countable additivity,
Fatou, dominated convergence, uniform integrability, plus exact counterexamples
when hypotheses are omitted. \\
Characteristic functions & Derivative-at-zero estimate, finite cyclic Fourier
expectation identity, and continuous effective inversion under explicit
Fubini/tail hypotheses. \\
Applications & Sample means, residual squares, moment determinacy on bounded
support, expected utility, indicators and Monte Carlo, center of mass, beta
balance, bounded Wald identity, variance, and finite-dimensional quantum
expectation. \\
\bottomrule
\end{longtable}

Historical claims, naming conventions, and etymology were not treated as
mathematical proof obligations.  The report's aim is not to reconstruct the
article's prose but to identify its formal mathematical core.

\subsection{Main constructive conclusion}

The expected-value operator is unusually well suited to \(\ipra\).  Its most
important algebraic laws are consequences of one primitive-recursive bounded
sum.  Even many probabilistic inequalities are pointwise finite inequalities
followed by that same sum.  The logical strength lies not in measure theory but
in the representation of the data.

The exact finite theory therefore needs no completed real line, sigma-algebra,
or countable additivity.  It already includes enough probability for
finite experiments, exact randomized algorithms, finite statistical models,
bounded stopping times, finite-dimensional quantum systems, and explicit
finite-error laws of large numbers.

For infinite objects, the correct constructive currency is a modulus.  An
absolute-tail modulus makes a countable expectation order independent.  An
\(L^1\) modulus makes expectations of simple approximants Cauchy.  A quadrature
modulus and an improper-tail modulus turn an integral into a regular real.  A
uniform-integrability modulus licenses limit interchange.  Once these data are
present, the proofs again reduce to the finite kernel.

\subsection{Corrections and cautions}

Three points deserve emphasis.

\begin{enumerate}
\item A conditionally convergent weighted series is not an expectation merely
  because one enumeration converges.  The outcomes have no canonical order;
  absolute convergence or a positive/negative-part definition is essential.
\item A symmetric principal value is not an expectation.  The Cauchy law has
  zero symmetric truncations but different asymmetric limits, while both
  absolute parts diverge.
\item The die/Chebyshev value \(8/15\) concerns the strict event
  \(1<X<6\), not the inclusive event that a die lies in its support.  The
  exact strict-event probability is \(2/3\), and the inclusive probability is
  one.
\end{enumerate}

\subsection{What the packet does not yet prove in Lean}

The human proofs are designed for compilation, but the packet stops before the
final kernel-checked endpoint.  In particular:
\begin{itemize}
\item the ambient Lean theorem-package records are not instantiated;
\item the natural decoders and common-denominator compiler are not implemented;
\item no concrete semantic checker has yet been compiled to \code{PR};
\item no new \code{Prf PRAAx} derivation is present;
\item the analytic modules are interfaces rather than completed interval
  arithmetic.
\end{itemize}
These limitations are explicit so that standard-model mathematics is not
mistaken for an internal proof.

\subsection{Further research questions}

\begin{enumerate}
\item \textbf{Minimal finite expectation kernel.}
  What is the smallest extension of the current ProveIt bounded-sum library
  that makes linearity, finite LOTUS, and Markov one- or two-line internal
  derivations?  A dedicated weighted-sum abstraction may reduce proof size
  substantially without increasing trusted strength.

\item \textbf{Canonical denominator compiler.}
  Can a verified compiler transform ambient \code{Rat} finite-sum theorems
  into signed-natural, denominator-cleared \code{PRAAx} goals automatically,
  while exposing a compact proof certificate checked by the existing kernel?

\item \textbf{Probability-specific reflection.}
  The Lagrange identity, variance expansion, and binomial expectation are
  polynomial identities after finite sums are normalized.  A reflection
  tactic that combines \code{MultiSums} normalization with \code{deriveNat}
  could cover a broad class of probability proofs.

\item \textbf{First-crossing library.}
  Kolmogorov's inequality, optional-stopping arguments, and bounded martingale
  theorems share a primitive-recursive first-hit construction.  A reusable
  first-crossing module may unlock many finitary stochastic-process results in
  \(\ipra\).

\item \textbf{Optimal effective hypotheses.}
  The analytic theorems here use convenient \(L^1\), tail, and probability
  moduli.  Which of these can be weakened while retaining a primitive-recursive
  output modulus?  In particular, can effective almost-sure convergence be
  represented more economically than convergence-in-probability plus uniform
  integrability?

\item \textbf{Constructive Radon--Nikodym fragments.}
  For which finitely generated algebras, computable compact spaces, or bounded
  density classes can a density be extracted in \(\ipra\) from a quantitative
  absolute-continuity modulus?  The general classical converse is too weakly
  specified, but important restricted forms may be fully effective.

\item \textbf{Exact cyclotomic Fourier layer.}
  Is a small verified cyclotomic-ring library simpler than regular sine/cosine
  names for finite Fourier inversion?  Such a layer would also benefit finite
  characteristic functions and discrete harmonic analysis.

\item \textbf{Formal expected utility and games.}
  The finite maximization theorem is immediate.  A natural next target is a
  verified constructive account of the problem of points, finite zero-sum
  games with rational mixed strategies, and backward induction, all phrased as
  expectation theorems.

\item \textbf{Moment problems at bounded complexity.}
  Finite support is exactly determined by finitely many moments.  What useful
  effective determinacy results hold for distributions with a supplied tail
  envelope, finite mixture models, or compact algebraic supports, within
  feasible \(\ipra\) proof size?

\item \textbf{Compression of analytic certificates.}
  The regular-real approach is transparent but potentially verbose.  Can
  Taylor, quadrature, and tail certificates be compressed into a single
  proof-producing evaluator whose correctness is established once in
  \(\mathrm{PRAAx}\)?
\end{enumerate}

\subsection{Final assessment}

A substantial portion of elementary expectation theory is not merely
constructively acceptable; it is primitive-recursively explicit.  The finite
kernel contains the conceptual heart of the subject, while the infinite theory
is a disciplined closure under named limits.  The remaining challenge is
proof engineering: implementing the encoders and turning the arguments already
reduced to finite algebra into compact checked \code{PRAAx} derivations.

\appendix
\section{Detailed algebraic certificates}\label{app:certificates}

\subsection{Denominator-cleared linearity}

Suppose \(X_i=a_i/D\) and \(Y_i=b_i/E\), where \(a_i,b_i\) are signed integer
codes and \(D,E>0\).  A common denominator is \(DE\), and
\[
 X_i+Y_i=\frac{Ea_i+Db_i}{DE}.
\]
The equality \(\E[X+Y]=\E[X]+\E[Y]\) becomes
\[
 \sum_iw_i(Ea_i+Db_i)
 =E\sum_iw_ia_i+D\sum_iw_ib_i.
\]
This is exactly \code{sum\_add} plus distributivity.  Notice that the total
weight \(W\) cancels from both sides before the object-language proof begins.

\subsection{Variance without field division}

Let \(A=\sum_iw_ia_i\), \(B=\sum_iw_ia_i^2\), and let the observable have
common denominator \(D\).  Then \(\mu=A/(WD)\).  Clearing the denominator
\(W^3D^2\) from
\(\operatorname{Var}(X)=\mathbb E[X^2]-\mu^2\) yields the economical
identity
\[
 \sum_iw_i(Wa_i-A)^2
 =W^2B-WA^2=W(WB-A^2).
\]
Indeed,
\begin{align*}
\sum_iw_i(W^2a_i^2-2Wa_iA+A^2)
&=W^2B-2WA\sum_iw_ia_i+A^2\sum_iw_i\\
&=W^2B-2WA^2+WA^2\\
&=W^2B-WA^2.
\end{align*}
Thus the denominator-cleared variance formula is a natural polynomial identity
using only \(A=\sum_iw_ia_i\) and \(\sum_iw_i=W\).

\subsection{Tail indicators}

For naturals \(x\le B\), define
\(T_B(x)=\sum_{k<B}\ind_{k<x}\).  The recurrences are
\[
 T_0(x)=0,
\]
and
\[
 T_{B+1}(x)=T_B(x)+\ind_{B<x}.
\]
An induction on \(B\) proves
\(T_B(x)=\min(x,B)\).  The successor step splits on the decidable comparison
\(B<x\).  Under \(x\le B\), this gives \(T_B(x)=x\).  This formulation avoids
an induction whose bound depends on the observable value and maps directly to
a two-variable open formula.

\subsection{Weighted Lagrange identity}

For completeness, expand the pairwise-square right side:
\begin{align*}
&\sum_{i<j}w_iw_j(u_iv_j-u_jv_i)^2\\
&=\sum_{i<j}w_iw_j(u_i^2v_j^2+u_j^2v_i^2-2u_iu_jv_iv_j)\\
&=\frac12\sum_{i\ne j}w_iw_j
  (u_i^2v_j^2-u_iu_jv_iv_j)\\
&=\frac12\left[
   \sum_{i,j}w_iw_j(u_i^2v_j^2-u_iu_jv_iv_j)
   -\sum_iw_i^2(u_i^2v_i^2-u_i^2v_i^2)
  \right]\\
&=\left(\sum_iw_iu_i^2\right)
  \left(\sum_jw_jv_j^2\right)
 -\left(\sum_iw_iu_iv_i\right)^2.
\end{align*}
For an arithmetic formalization one should avoid the factor \(1/2\): split the
full off-diagonal sum into the two ordered triangles and prove directly that
they are equal under index swap.

\subsection{Finite geometric derivative identity}

For \(q\ne1\), the finite geometric sum is
\[
 G_N(q)=\sum_{k=0}^{N-1}q^k,
 \qquad (1-q)G_N(q)=1-q^N.
\]
The weighted version needed for the geometric mean is proved without formal
differentiation.  Let
\(S_N=\sum_{k=1}^{N}kq^{k-1}\).  Then
\begin{align*}
(1-q)S_N
 &=1+q+\cdots+q^{N-1}-Nq^N\\
 &=G_N(q)-Nq^N.
\end{align*}
Multiplying again by \(1-q\) gives
\[
 (1-q)^2S_N=1-(N+1)q^N+Nq^{N+1}.
\]
Both equations follow by aligning two finite sums and canceling equal terms;
they are ideal quantifier-free induction lemmas.

\subsection{First-crossing indicators}

For a finite path \(S_1,\ldots,S_n\), define
\[
 C_k=\ind_{\abs{S_k}\ge a}
      \prod_{j<k}\ind_{\abs{S_j}<a}.
\]
Products of Boolean indicators are natural \(0/1\) terms.  Then
\[
 C_kC_l=0\quad(k\ne l),
\]
because the later first-crossing indicator contains a factor asserting that
the earlier threshold was not crossed.  Furthermore,
\[
 \ind_{\max_{j\le n}\abs{S_j}\ge a}=\sum_{k=1}^{n}C_k.
\]
The right side is at most one by disjointness and is one exactly when bounded
search finds a least crossing.  This explicit formula can replace an abstract
finite-event partition in the Kolmogorov proof.

\section{Validation and reproducibility}\label{app:validation}

\subsection{Exact regression program}

The file \code{code/validate\_expected\_value.py} uses Python's
\code{fractions.Fraction} exclusively for pass/fail arithmetic.  The recorded
run in \code{validation/validation.json} reports:

\begin{center}
\begin{tabular}{@{}lr@{}}
\toprule
Test group & Exact assertions \\
\midrule
Article examples & 3 \\
Finite normalization, bounds, parts, indicators, variance & 24,220 \\
Linearity and monotonicity & 37,462 \\
Tail formulas & 12,784 \\
Markov, Chebyshev, and die audit & 3,748 \\
Cauchy--Schwarz and Lagrange identity & 63,750 \\
Independence and dependence examples & 7 \\
Finite conditional expectation & 1 \\
Binomial and sample-mean formulas & 36 \\
Geometric and Poisson identities & 110 \\
Divergence and rearrangement diagnostics & 87 \\
Convergence and sigma-additivity counterexamples & 21 \\
Finite moment determinacy & 3,300 \\
Center of mass & 2,744 \\
Kolmogorov maximal inequality & 44 \\
\midrule
Total & 148,317 \\
\bottomrule
\end{tabular}
\end{center}

The tests enumerate many small laws and observables, verify the algebraic
identities exactly, and exercise counterexamples.  They do not prove universal
quantification, and the Python interpreter is not part of the trusted theory.
Their role is to detect transcription errors before Lean formalization.

\subsection{Reproduction commands}

From the archive root:

\begin{verbatim}
python3 code/validate_expected_value.py \
  --output validation/validation.json
cd paper
latexmk -pdf -interaction=nonstopmode -halt-on-error \
  expected_value_ipra.tex
\end{verbatim}

The source tree records the consulted Wikipedia revision and the ProveIt commit
in \code{provenance/SOURCES.md}.  The top-level \code{MANIFEST.txt} lists the
intended archive contents.

\section{A compact object-language checklist}\label{app:checklist}

For each theorem file, a reviewer should be able to answer ``yes'' to all of
the following.

\begin{enumerate}
\item Are all mathematical inputs natural-number codes with explicit validity
  predicates?
\item Is every rational denominator proved positive before cross
  multiplication or cancellation?
\item Are all finite arrays accessed through proved bounded decoders?
\item Is each bounded sum/product recurrence represented by a named \code{PR}
  term with proved defining equations?
\item Does the theorem checker reject a mathematical counterexample rather
  than return zero on every input?
\item Is the open goal quantifier-free, so that every induction use is an
  instance of the permitted rule?
\item Are bounded universal conditions compiled to finite failure sums or an
  equivalent quantifier-free code?
\item Is the semantic correctness theorem separate from the internal
  derivation?
\item Is universal closure performed only after the open derivation is
  complete?
\item Is transport to the native calculus performed only through the existing
  proved equivalence?
\end{enumerate}

Meeting this checklist would turn the present blueprint into a genuinely
kernel-checked \(\ipra\) development.


\begin{thebibliography}{9}
\bibitem{wikipediaEV}
Wikipedia contributors, ``Expected value,'' revision consulted 16 November 2026.

\bibitem{proveit}
V. Reshetnikov, \emph{ProveIt}, repository revision
\code{df3ff8cc28be} (full hash in provenance), 16 November 2026.

\bibitem{conciseipra}
V. Reshetnikov, ``Concise presentations of intuitionistic primitive recursive arithmetic,''
ProveIt/Logic/PrimitiveRecursiveArithmetic/ConciseIPRA.

\bibitem{cltpacket}
``Constructive central limit theorems in iPRA,'' ProveIt proof packet,
15 November 2026.
\end{thebibliography}

\end{document}
